\providecommand{\dirfiles}[1]{\par\smallskip\noindent
  {\footnotesize\raggedright\emph{Files:} #1.\par}\smallskip}
\providecommand{\dirP}{\mathbb{P}}
\providecommand{\dirE}{\mathbb{E}}
\providecommand{\dirU}{\mathrm{U}}
\section{CM coefficient fields, reflected fields and uniform Galois bounds}\label{sec:cmf}

\subsection{Coefficient fields of CM newforms}\label{dir:sec:cm}
\sources{grh-lean-transfers}

A normalized newform $f=\sum_{n\ge1}a_nq^n$ has a coefficient field
$K_f=\Q(a_1,a_2,\dots)$. Fixing the weight and the degree of this field still allows
the level to grow without bound. Martin's Conjecture~1(a)
\cite{Martin2025Rationality} asks whether only finitely many fields occur at fixed
weight and degree for trivial nebentypus. For the forms with complex multiplication
(CM) the question can be decided, and the decisive zeros belong to quadratic Dirichlet
$L$-functions; no hypothesis on the $L$-functions of the modular forms is needed. In
weight two the classifications below turn out to be unconditional.

Throughout, $f$ is a CM newform of weight $k\ge2$, $K=K_f$, $d=[K:\Q]$, $w=k-1$, and
$E$ is the imaginary quadratic CM field, of fundamental discriminant $D$ and class
number $h_E$. The CM correspondence realizes $f$ as the theta series of an algebraic
Hecke character $\psi$ of $E$ of infinity type $(w,0)$ and conductor $\mathfrak m$:
$a_n$ is the sum of $\psi(\mathfrak a)$ over the integral ideals of norm $n$, with
$\psi$ extended by zero to ideals meeting the conductor. Let $L$ be the field generated
by the values of $\psi$, and let $\eta\colon(\OO_E/\mathfrak m)^\times\to\C^\times$ be
the finite character of $\psi$, so that
\begin{equation}\label{dir:eq:eta}
 \psi((a))=\eta(a)\,a^w\qquad(a\in E^\times\text{ a unit at }\mathfrak m).
\end{equation}

\begin{lemma}[The value field; \stU]\label{dir:lem:valuefield}
For every nebentypus, $L=EK$. In particular $[L:E]=d/[K\cap E:\Q]$ divides $d$.
\end{lemma}

\begin{proof}
Clearly $K\subseteq L$. For a prime $p=\mathfrak p\bar{\mathfrak p}$ that splits in $E$
and is prime to the conductor, put $u=\psi(\mathfrak p)$ and
$v=\psi(\bar{\mathfrak p})$. The coefficients at $p$ and $p^2$ give
\[
 u+v=a_p,\qquad uv=a_p^2-a_{p^2}.
\]
Let $h$ be the exponent of the ray class group modulo $\mathfrak m$. Then
$\mathfrak p^h=(\alpha)$ and $\bar{\mathfrak p}^h=(\beta)$ with generators congruent to
$1$ modulo $\mathfrak m$, so $u^h=\alpha^w$ and $v^h=\beta^w$ belong to $E$. Their
principal ideals are distinct, and neither is invariant under conjugation.
Consequently $u^h$ generates $E$ over $\Q$, so $E\subseteq L$. Every automorphism of
$\overline{\Q}$ fixing $EK$ preserves the pair of roots $u,v$ of
$X^2-a_pX+(a_p^2-a_{p^2})$ and fixes their distinct $h$th powers; it therefore fixes
each root, and $u,v\in EK$. At a split prime where only one factor is prime to
$\mathfrak m$, its value equals $a_p$. At an inert prime the sole prime ideal has value
$a_{p^2}$, and at a ramified prime its value is $a_p$, whenever it is prime to the
conductor. Prime ideals generate the ideal group, which proves $L=EK$.
\end{proof}

\begin{lemma}[The class-group exponent; \stU]\label{dir:lem:cmexponent}
The exponent $e_E$ of $\Cl(E)$ divides $w[L:E]$, hence divides $(k-1)d$.
\end{lemma}

\begin{proof}
For an ideal $\mathfrak a$ prime to $\mathfrak m$ write $\mathfrak a^h=(\gamma)$ with
$\gamma\equiv1\pmod{\mathfrak m}$. Taking $h$th powers of ideals and using
$\psi(\mathfrak a)^h=\gamma^w$ gives $(\psi(\mathfrak a))=(\mathfrak a\OO_L)^w$ in $L$.
The ideal norm from $L$ to $E$ shows that $\mathfrak a^{w[L:E]}$ is principal. Every
ideal class has such a representative.
\end{proof}

\begin{lemma}[Odd degree forces trivial nebentypus; \stU]\label{dir:lem:oddneben}
If $d$ is odd, then $f$ has trivial nebentypus and $K$ is totally real.
\end{lemma}

\begin{proof}
A field of odd degree has a real embedding. This embedding conjugates $f$ to a newform
with real coefficients, so Sch\"utt's Corollary~1.2(i) \cite{Schutt2005CM} forces its
nebentypus, and therefore the original one, to be trivial. For trivial nebentypus the
coefficient field is totally real.
\end{proof}

\begin{theorem}[Finiteness with varying character; \stG]\label{dir:thm:cmfinite}
Fix integers $k\ge2$ and $d\ge1$. Only finitely many isomorphism classes of fields
$K_f$ occur among the CM newforms of weight $k$ with $[K_f:\Q]=d$, as both the level
and the nebentypus vary.
\end{theorem}

\begin{proof}
Martin proves the case of trivial character under quadratic Dirichlet GRH
\cite[Theorem~1]{Martin2025Rationality}, and his Remark~1.1 anticipates a
fixed-character analogue. The class-group and value-field arguments below are his; the
identity $L=EK$ of Lemma~\ref{dir:lem:valuefield} is what allows the character to vary.

\emph{Bounding the CM discriminant.} If $p=\mathfrak p\bar{\mathfrak p}$ splits in $E$,
take $\mathfrak p^{e_E}=(\alpha)$. Unequal valuations at the two primes above $p$ force
$\alpha\notin\Q$. Writing $\alpha=(a+b\sqrt D)/2$ with $b\ne0$ gives
\[
 p^{e_E}=\abs{\Nm_{E/\Q}\alpha}\ge\abs D/4 .
\]
Put $T=\log\abs D$ and $M=(k-1)d$, so that $e_E\mid M$ by Lemma~\ref{dir:lem:cmexponent}.
If $T>25$, Theorem~3 of Elsenhans, Kl\"uners and Nicolae
\cite{Elsenhans2020Imaginary} gives a split prime
$p\le(1.881T+6.18)^2$ under the Riemann hypothesis for $\zeta_E(s)=\zeta(s)L(s,\chi_D)$,
which \cref{src:thm:dgrh} supplies. Hence
\begin{equation}\label{dir:eq:cmcutoff}
 T\le\log4+2M\log(1.881T+6.18).
\end{equation}
The right side grows only logarithmically in $T$, so \eqref{dir:eq:cmcutoff} is an
explicit discriminant cutoff for every weight and coefficient degree; only finitely
many fields $E$ are possible.

\emph{Finishing the field count.} For any one of these $E$, Martin's Corollary~3.6
gives finitely many value fields $L/E$ of each bounded degree, uniformly in the
conductor. Indeed the image of $\psi$ on principal ideals generates $E(\mu_r)$, where
$r$ is the order of $\eta$; thus $\varphi(r)\le[L:\Q]\le2d$. Choose ideal-class
generators $\mathfrak t_i$ with $\mathfrak t_i^{n_i}=(\theta_i)$. Martin's
Proposition~3.4 places $L$ in the field obtained from $E$ by adjoining $\mu_{re_E}$ and
the $n_i$th roots of $\theta_i^w$, independently of $\mathfrak m$. There are finitely
many possible $r$, and each resulting number field has finitely many subfields.
\end{proof}

For example, in weight two and degree seven the class-group exponent divides seven,
and Table~2 of \cite{Elsenhans2020Imaginary} yields
$\abs D\le3.9\cdot10^{30}$. In weight two there is an independent unconditional
reduction: the divisibility
\begin{equation}\label{dir:eq:yang}
 h_E\mid[L:E]\mid d\qquad(k=2)
\end{equation}
of Yang \cite[Theorem~3.1]{Yang2003CM}, applied to modular forms by Gonz\'alez
\cite[Section~3]{Gonzalez2011Endomorphism}, replaces the class-group exponent by the
class number, and \cite[Theorem~3.1]{Gonzalez2011Endomorphism} already proves
weight-two finiteness in all degrees and for all characters. The interest below is in
the exact lists over all levels.

\paragraph{The genus field.}
After Theorem~2 of \cite{Martin2025Rationality}, Martin observes from his tables that
when a weight-two CM newform of trivial nebentypus has quadratic coefficient field $K$
and $h_E=2$, the field $EK$ is the Hilbert class field of $E$, and he points out the
lack of a non-computational proof. The following containment explains the equality,
extends to every even weight, and removes the analytic hypothesis from his
classification in coefficient degree two.

\begin{theorem}[The genus field inside the value field; \stU]\label{dir:thm:genus}
Let $f$ be a normalized CM newform of trivial nebentypus and even weight $k\ge2$. Let
$t$ be the number of distinct primes dividing $D$, and $G_E$ the genus field of $E$.
Then
\[
 G_E\subseteq L=EK,\qquad G_E\cap\R\subseteq K,\qquad 2^{t-1}\mid[K:\Q].
\]
In particular, if $h_E=[K:\Q]=2$, then $EK$ is the Hilbert class field $H_E$ of $E$,
and $K$ has discriminant equal to the positive signed prime-discriminant factor of $D$.
\end{theorem}

\begin{proof}
Here $w=k-1$ is odd. By the formula for the nebentypus of the theta series of $\psi$,
the restriction of $\eta$ to rational units is the primitive quadratic character
$\chi_D$ away from the moduli. The finite image $U$ of $\eta$ lies in $L$ by
\eqref{dir:eq:eta}, since $E\subseteq L$ and each residue unit has a lift. Group the
local components of $\eta$ over each rational prime $v$ into $\eta_v$. Their
restrictions to diagonal rational units are exactly the corresponding local components
of $\chi_D$; this follows by factoring the equality of Dirichlet characters into
prime-power components using the Chinese remainder theorem, and it allows $\mathfrak m$
to be asymmetric over a split prime.

Fix an odd prime $s\mid D$ and let $P_s^2=(s)$ be its ramified prime ideal. By weak
approximation choose $x\in E^\times$ such that $\mathfrak a=xP_s$ is prime to
$\mathfrak m$. Then $\mathfrak a^2=(\theta)$ for $\theta=sx^2$, and \eqref{dir:eq:eta}
gives
\begin{equation}\label{dir:eq:genussquare}
 b^2=s\,\eta(\theta),\qquad b=\frac{\psi(\mathfrak a)}{x^ws^{(w-1)/2}}\in L .
\end{equation}
We compute $\eta(\theta)$ in $U/U^2$. At a prime $v\ne s$ in the conductor, $\theta$
equals $s$ times a local unit square at each relevant place. At $s$ use instead
\[
 \theta=(D/s)\,(sx/\sqrt D)^2 ,
\]
in which $D/s$ is a rational local unit and the second factor is a local unit square.
Put $e_s=(-1)^{(s-1)/2}$, $s^*=e_ss$ and $d'=D/s^*$. The factorization into primitive
quadratic characters gives $\chi_D=\chi_{s^*}\chi_{d'}$ on common units, with coprime
conductors. Away from $s$ the rational local components evaluated at $s$ multiply to
$\chi_{d'}(s)=\legendre{d'}{s}$. At $s$ the restriction to rational units is
$\legendre{\cdot}{s}$. Multiplying in $U/U^2$ yields
\[
 \eta(\theta)\equiv\legendre{d'}{s}\legendre{D/s}{s}
 =\legendre{e_s(d')^2}{s}=e_s\pmod{U^2}.
\]
All signs belong to $U$ since $\eta(-1)=-1$. Consequently there is $u\in U\subseteq L$
with $\eta(\theta)=e_su^2$, and \eqref{dir:eq:genussquare} gives $(b/u)^2=s^*$. Thus
$\sqrt{s^*}\in L$ for every odd prime $s\mid D$, with no restriction on the conductor,
the order of $\eta$ or the class of $P_s$. The square-root identity at a ramified prime
has a weight-two antecedent in \cite[Lemma~3.3]{Yang2003CM}.

Gauss's genus theory gives $G_E=\Q(\sqrt{d_1},\dots,\sqrt{d_t})$, where $D=\prod_jd_j$
is the factorization into signed prime discriminants, and $[G_E:E]=2^{t-1}$;
equivalently $\Gal(G_E/E)$ is the class group modulo squares \cite{LemmermeyerGenus}.
The odd factors are the $s^*$ just obtained. If a dyadic factor $-4$, $8$ or $-8$
occurs, its square root also belongs to $L$, by dividing $\sqrt D\in E$ by the product
of the square roots of the odd factors. This proves $G_E\subseteq L$. Martin's
Lemma~2.2 identifies $L=EK$ (compare Lemma~\ref{dir:lem:valuefield}), and $K$ is totally
real. Thus $L\cap\R=K$ and $[L:E]=[K:\Q]$, which gives the two remaining conclusions.
For $h_E=2$ the class group has trivial subgroup of squares, so $G_E=H_E$ has degree
two over $E$; when $[K:\Q]=2$ the containment is therefore an equality. The two signed
prime discriminants have opposite signs, and the positive one defines the unique real
quadratic subfield.
\end{proof}

\begin{corollary}[Quadratic coefficient fields in weight two; \stU]
\label{dir:cor:cmquadratic}
The real quadratic fields that occur as coefficient fields of weight-two CM newforms
of trivial nebentypus are exactly those with the $23$ discriminants
\[
\begin{gathered}
 5,8,12,13,17,21,24,28,29,33,37,41,\\
 44,57,61,76,89,129,172,201,268,489,652 .
\end{gathered}
\]
\end{corollary}

\begin{proof}
Suppose $k=2$ and $[K:\Q]=2$. Martin's Corollary~3.2 implies $\exp\Cl(E)\mid2$, so the
class group equals its quotient by squares, and by Theorem~\ref{dir:thm:genus} its order
$2^{t-1}$ divides two. Thus $h_E\in\set{1,2}$. The unconditional classifications of the
imaginary quadratic fields of class numbers one and two \cite{Watkins2004Class} provide
the base fields to check, and Martin's local arguments and constructions for these
fields give the list, which is his Table~1.1 with the existence proofs of his
Sections~5--7. The finite positive constructions and the specific class-number-one
exclusions do not use ERH.
\end{proof}

In \cite{Martin2025Rationality} this list is the degree-two part of Theorem~2, proved
there under quadratic Dirichlet GRH. Earlier work gives another route in weight two:
\eqref{dir:eq:yang} gives $h_E\mid[K:\Q]$, and \cite[Theorem~3.2]{Gonzalez2011Endomorphism}
classifies the $83$ pairs of CM field and quadratic coefficient field when the
nebentypus is allowed to vary. The genus argument is an alternative proof; its local
sign uses the trivial nebentypus and works in every even weight.

\paragraph{Degree three.}

\begin{theorem}[Cubic coefficient fields; \stU]\label{dir:thm:cmcubic}
In weight two exactly eighteen cubic fields occur as coefficient fields of CM
newforms, across all levels and nebentypus. Their discriminants are
\[
\begin{gathered}
49,\ 81,\ 321,\ 621,\ 837,\ 993,\ 1593,\ 1929,\ 2241,\ 3753,\\
5697,\ 7641,\ 8289,\ 10233,\ 13473,\ 14769,\ 23841,\ 24489 ,
\end{gathered}
\]
and each field is determined up to isomorphism by its discriminant.
\end{theorem}

\begin{proof}
By Lemma~\ref{dir:lem:oddneben} the nebentypus is trivial, and by
Lemmas~\ref{dir:lem:valuefield} and~\ref{dir:lem:cmexponent} we have $[L:E]=3$ and
$\exp\Cl(E)\mid3$.

Assume $h_E>1$. We first show that $\eta$ has order two. If its order is $r$, then
$\Q(\zeta_r)\subseteq L$, so $\varphi(r)\mid6$; also $r$ is even since $\eta(-1)=-1$.
If $\varphi(r)=2$, the cyclotomic quadratic field would have to equal $E$, as $[L:E]$
is odd, but $\Q(i)$ and $\Q(\sqrt{-3})$ have class number one. If $\varphi(r)=6$ then
$r=14$ or $18$; the sextic cyclotomic field would equal $L$, whose quadratic subfield
would be $\Q(\sqrt{-7})$ or $\Q(\sqrt{-3})$, again of class number one. Only $r=2$
remains.

For an ideal class $c$ choose a fractional ideal $\mathfrak a$ representing $c$ and
prime to the conductor, and choose $\theta_c\in E^\times$ with
$\mathfrak a^3=(\theta_c)$. The assignment
\[
 \Cl(E)\longrightarrow E^\times/(E^\times)^3,\qquad c\longmapsto[\theta_c]
\]
is a well-defined injective homomorphism: a change of representative multiplies
$\theta_c$ by a cube times a unit, and both units $\pm1$ are cubes; the homomorphism
law follows from ideal multiplication; and if $\theta_c=b^3$, unique factorization of
fractional ideals gives $\mathfrak a=(b)$. Every $\theta_c$ acquires a cube root in
$L$, because $\psi(\mathfrak a)^3=\eta(\theta_c)\theta_c=\pm\theta_c$. Set
$F=E(\zeta_3)$. Since $[F:E]=2$, the norm from $F$ to $E$ shows that
$E^\times/(E^\times)^3\to F^\times/(F^\times)^3$ is injective: if $b$ becomes a cube,
then $b^2$ is a cube in $E$, so $b$ is a cube already. By the Kummer correspondence
over $F$, adjoining cube roots of a subgroup of $F^\times/(F^\times)^3$ of order $3^m$
gives an extension of degree $3^m$. The subgroup just constructed has order $h_E$ and
its radicals all lie in $FL$. It follows that
\[
 h_E\mid[FL:F]\mid3 .
\]
In particular every nontrivial class group occurring here is cyclic of order three;
this excludes fields with arbitrarily many independent classes of order three without
enumerating them. (The divisibility \eqref{dir:eq:yang} also gives $h_E\mid3$
directly.)

There are sixteen imaginary quadratic fields of class number three, with largest
absolute discriminant $907$; the unconditional classification goes back to Oesterl\'e
and is included in \cite{Watkins2004Class}. For these sixteen fields Martin's
Corollary~4.5 and Section~6.2 give the unique cubic coefficient field in each case,
his Table~6.3 gives its discriminant, and the constructions there establish
occurrence. For class number one his Section~5.3 gives just $\Q(\zeta_7)^+$ and
$\Q(\zeta_9)^+$. These are the eighteen entries displayed, and by
\cite[Table~1.2]{Martin2025Rationality} each is identified by its discriminant.
\end{proof}

Martin's proof of the cubic part of his Theorem~2 used a GRH-conditional list of the
fields of class-group exponent three; the change here is the uniform exclusion of class
groups of order greater than three.

\paragraph{Degree five.}
Martin singles out degree five as the next case \cite[Remark~1.2]{Martin2025Rationality}.

\begin{theorem}[Quintic coefficient fields; \stU]\label{dir:thm:cmquintic}
Across all levels and nebentypus, exactly twenty-six quintic fields occur as
coefficient fields of weight-two CM newforms. One is the maximal real subfield
$\Q(\zeta_{11})^+$. Each of the other twenty-five comes from a different imaginary
quadratic field of class number five; for $D=-q$ it is generated by a root of
\begin{equation}\label{dir:eq:quinticpoly}
 P_{p,a}(X)=X^5-5pX^3+5p^2X-a ,
\end{equation}
with $(q,p,a)$ as in \cref{dir:tab:quintic}.
\end{theorem}

\begin{proof}
\emph{Reduction.} By Lemma~\ref{dir:lem:oddneben} the nebentypus is trivial. Put $L=EK$, so
that $[L:E]=5$ and $\exp\Cl(E)\mid5$. For $h_E>1$, $E$ has only the units $\pm1$. Let
$r$ be the order of $\eta$. As $\varphi(r)\mid10$ and $r$ is even, $r$ is one of
$2,4,6,22$. Orders four and six would force $E$ to be the quadratic cyclotomic field
inside $L$, since $[L:E]$ is odd, and order twenty-two would force $L=\Q(\zeta_{11})$
and hence $E=\Q(\sqrt{-11})$. All three exceptional quadratic fields have class number
one. Thus $r=2$ when $h_E>1$. For each class $[\mathfrak a]$ use a representative prime
to the conductor and write $\mathfrak a^5=(\theta)$. As in the proof of
Theorem~\ref{dir:thm:cmcubic}, the classes $[\theta]$ form a subgroup of
$E^\times/(E^\times)^5$ isomorphic to $\Cl(E)$, and each $\theta$ has a fifth root in
$L$ since $\psi(\mathfrak a)^5=\pm\theta$. These classes stay independent over
$F=E(\zeta_5)$, as $[F:E]$ divides four. Kummer theory gives $h_E\mid[FL:F]=5$. Only
$h_E=1$ and $h_E=5$ remain.

\emph{Class number one.} Principal ideals supply all values, so $L=E(\mu_r)$. For this
field to have degree ten over $\Q$ the cyclotomic argument leaves only
$E=\Q(\sqrt{-11})$ and $r=22$. To realize it, reduce modulo $11\OO_E$. The quotient is
the ring of dual numbers over $\F_{11}$, with unit group
$\F_{11}^\times\times(1+\epsilon\F_{11})$, $\epsilon^2=0$. Choose the Legendre
character on the first factor and a character of order eleven on the second. The
resulting $\eta$ sends $-1$ to $-1$ and restricts on rational units to $\chi_{-11}$.
Hence $\psi((\alpha))=\alpha\eta(\alpha)$ defines a Hecke character of type $(1,0)$,
primitive because the part of order eleven is nontrivial. Its newform has trivial
nebentypus and coefficient field $\Q(\zeta_{11})^+$.

\emph{Class number five: uniqueness.} Genus theory and the odd class number give a
single prime discriminant, necessarily $D=-q$ with $q\equiv3\pmod4$ an odd prime; the
purely dyadic discriminants have class number one. Let $\mathfrak t$ be an integral
ideal representing a generator of $\Cl(E)$, prime to $q$, of norm $n$, and choose
$\theta$ with $\mathfrak t^5=(\theta)$. For any conductor take a representative
$\alpha\mathfrak t$ coprime to it. Its character value, divided by $\alpha$, is a
fifth root of $\pm\theta$. The values on principal ideals generate $E$, and this one
value generates the rest of $L$ over $E$; the sign can be absorbed in the choice of
fifth root. The polynomial $X^5-\theta$ is irreducible over $E$: if a root generated
an extension of degree $1\le j<5$, its norm would show that $\theta^j$ is a fifth
power, forcing $\theta$ to be a fifth power, contrary to the order of $\mathfrak t$.
Thus all choices yield isomorphic $L$ over $E$. As $L$ is a totally imaginary quadratic
extension of its totally real coefficient field, its maximal totally real subfield
recovers $K_f$. Hence there is at most one coefficient field for each $E$.

\emph{Class number five: existence.} Choose as conductor the prime above $q$ and as
$\eta$ the Legendre character on its residue field. It sends $-1$ to $-1$ and its
rational restriction is $\chi_D$. The rule $\psi((\alpha))=\eta(\alpha)\alpha$ on
principal ideals extends to all ideals prime to the conductor by assigning to
$\mathfrak t$ any fifth root of $\eta(\theta)\theta$: every ideal is a principal ideal
times a power of $\mathfrak t$, and the only relation of $\mathfrak t$ modulo principal
ideals is the fifth-power relation used here. This character is primitive at $q$; its
newform has weight two, trivial nebentypus and coefficient degree five. If
$z=\psi(\mathfrak t)$ and $\Theta=z^5$, then $\bar z=n/z$ because $\abs\theta^2=n^5$.
The coefficient field is generated by $x=z+n/z$, a root of
\[
 X^5-5nX^3+5n^2X-\Tr_{E/\Q}(\Theta).
\]
Indeed $x$ is real and $[E(z):E]=5$, while $z$ satisfies $Z^2-xZ+n=0$, so
$x\notin\Q$.

\emph{Different CM fields give different quintic fields.} Set $B=E(\zeta_5)$ and
$M=B(z)$. Since $[B:E]$ divides four, $\Theta$ remains a non-fifth-power in $B$, so
$M/B$ is cyclic of degree five. Conjugation takes $\Theta$ to $n^5/\Theta$, so $M$ is
normal over $\Q$. It is unramified outside $5q$: at all other places $B$ is unramified
and the valuation of $\Theta$ is a multiple of five, so that after division by a fifth
power of a uniformizer one adjoins a fifth root of a unit, whose reduction is separable
at that residue characteristic. At $q$, an inertia element restricting nontrivially to
$E$ fixes $\zeta_5$ and sends $z$ to $\zeta_5^jn/z$ for some $j$. The five conjugates
of $x$ are $x_i=\zeta_5^iz+\zeta_5^{-i}n/z$. They are distinct, since an equality for
two different $i$ would make $z^2/n$ a root of unity, so that $\mathfrak t^5$ and
$\bar{\mathfrak t}^5$ would be equal ideals and $[\mathfrak t]=[\mathfrak t]^{-1}$. The
inertia element acts on the five conjugates by $i\mapsto-i-j$ and therefore acts
nontrivially. This forces ramification at $q$ in the quintic field: if it were
unramified at $q$, the compositum of all its conjugates would be unramified there as
well. Thus $q$ is the unique ramified rational prime other than five, and it recovers
the CM field. The field $\Q(\zeta_{11})^+$ ramifies only at eleven, and
$\Q(\sqrt{-11})$ has class number one, so it is distinct from every class-number-five
example.

\emph{The count and the polynomials.} The unconditional classification of class number
five contains twenty-five fields, with largest absolute discriminant $2683$
\cite{Watkins2004Class}. In \cref{dir:tab:quintic}, $p$ is the norm of a nonprincipal
split prime ideal and $4p^5=a^2+qb^2$. The integer $(a+b\sqrt{-q})/2$ generates the
fifth power of one of the two prime ideals over $p$: it has the required norm, and
$p\nmid a$, so its two conjugates cannot both be divisible by the same prime over $p$.
Taking $\mathfrak t$ to be this prime gives $n=p$ and $\Tr\Theta=\pm a$; changing the
sign of the trace sends the root to its negative and gives the same field.
\end{proof}

\begin{table}[ht]
\centering\small
\begin{tabular}{rrrrr@{\qquad}rrrrr}
\toprule
$q$ & $p$ & $a$ & $b$ & $v_5$ & $q$ & $p$ & $a$ & $b$ & $v_5$\\
\midrule
47 & 2 & 9 & 1 & 5 & 691 & 5 & 38 & 4 & 5\\
79 & 2 & 7 & 1 & 5 & 739 & 5 & 26 & 4 & 3\\
103 & 2 & 5 & 1 & 5 & 787 & 7 & 59 & 9 & 5\\
127 & 2 & 1 & 1 & 5 & 947 & 3 & 5 & 1 & 5\\
131 & 3 & 29 & 1 & 5 & 1051 & 5 & 107 & 1 & 3\\
179 & 3 & 16 & 2 & 5 & 1123 & 7 & 239 & 3 & 3\\
227 & 3 & 8 & 2 & 5 & 1723 & 11 & 149 & 19 & 3\\
347 & 3 & 25 & 1 & 3 & 1747 & 17 & 1565 & 43 & 5\\
443 & 3 & 23 & 1 & 3 & 1867 & 7 & 4 & 6 & 5\\
523 & 7 & 220 & 6 & 5 & 2203 & 7 & 255 & 1 & 5\\
571 & 5 & 58 & 4 & 5 & 2347 & 17 & 1810 & 32 & 5\\
619 & 5 & 109 & 1 & 5 & 2683 & 11 & 344 & 14 & 5\\
683 & 3 & 17 & 1 & 5 &  & & & & \\
\bottomrule
\end{tabular}
\caption{The twenty-five quintic CM coefficient fields attached to class number five:
$D=-q$, $4p^5=a^2+qb^2$, defining polynomial $P_{p,a}$ of \eqref{dir:eq:quinticpoly}.
The column $v_5$ is the $5$-adic valuation of the field discriminant $q^2\,5^{v_5}$
(Corollary~\ref{dir:cor:quinticdisc}).}\label{dir:tab:quintic}
\end{table}

The remaining field $\Q(\zeta_{11})^+$ has defining polynomial
$X^5+X^4-4X^3-3X^2+3X+1$, obtained by writing $1+\zeta+\dots+\zeta^{10}=0$ in terms of
$X=\zeta+\zeta^{-1}$. The first row of the table gives $X^5-10X^3+20X-9$ for $D=-47$.
The accompanying program counts reduced primitive binary quadratic forms up to
Watkins's discriminant bound, finds the displayed norm identities and checks every
polynomial modulo primes. For each of the $26$ polynomials the archive gives a prime of
irreducible reduction, and for each of the $\binom{26}{2}=325$ pairs a prime where both
reductions are squarefree and the factor degrees disagree. A
squarefree reduction ensures that the prime does not divide the polynomial
discriminant, hence not the index of the order; its factor degrees are then the
residue degrees in the field, and unequal patterns certify nonisomorphism. The
canonical form of level $q^2$ and coefficient degree $h(-q)$ belongs to earlier work
\cite{GonzalezLario2008Qcurves,GrossHecke}; what is added in
Theorem~\ref{dir:thm:cmquintic} is exhaustiveness over all levels, the distinction by the
ramified prime and the explicit polynomial list.
\dirfiles{\artifact{grh-lean-transfers/checks/quintic_fields.py},
\artifact{grh-lean-transfers/checks/quintic_irreducible.csv},
\artifact{grh-lean-transfers/checks/quintic_witnesses.csv}}

\paragraph{Every odd prime degree at least seven.}
Define the Dickson polynomials by $D_0(X,p)=2$, $D_1(X,p)=X$ and
$D_n=XD_{n-1}-pD_{n-2}$ for $n\ge2$.

\begin{theorem}[Prime degree at least seven; \stU]\label{dir:thm:cmoddprime}
For each prime $\ell\ge7$, the isomorphism classes of degree-$\ell$ coefficient fields
of weight-two CM newforms, of any level and nebentypus, are in bijection with the
imaginary quadratic fields $E$ of class number $\ell$. The field attached to $E$ of
discriminant $-q$ is the real field generated by a root of
\[
 D_\ell(X,p)-\Tr_{E/\Q}\alpha ,
\]
where $P$ is a nonprincipal prime ideal of $E$ of degree one, of norm $p\nmid q\ell$,
and $P^\ell=(\alpha)$. In particular there are $31$ such fields in degree seven, $41$
in degree eleven and $37$ in degree thirteen.
\end{theorem}

\begin{proof}
By Lemma~\ref{dir:lem:oddneben} the nebentypus is trivial; let $L=EK$. The divisibility
\eqref{dir:eq:yang} gives $h_E\mid[L:E]=\ell$.

\emph{Restricting the finite character.} If $r$ is the order of $\eta$, then $r$ is
even and $\varphi(r)\mid2\ell$ because $\mu_r\subseteq L$. The possibilities are
$r=2,4,6$, and $r=2(2\ell+1)$ if $2\ell+1$ is prime. Indeed $\varphi(r)=2\ell$ requires
an odd prime factor of $r$ whose predecessor divides $2\ell$ and is even; that prime is
$3$ or $2\ell+1$, and a power of $3$ cannot supply $\ell\ge7$ in the totient.

In the last case $q=2\ell+1$ and $L=\Q(\zeta_q)$, so its unique quadratic subfield is
$E=\Q(\sqrt{-q})$. This field cannot have class number one: none of the prime absolute
discriminants $3,7,11,19,43,67,163$ of class number one has $(q-1)/2$ prime and at
least seven. It cannot have class number $\ell$ either. For any negative prime
discriminant $-q$ with $q\ge7$, a reduced form $(a,b,c)$ satisfies $a\le\sqrt{q/3}$;
there are at most $a$ odd values $-a<b\le a$, by the boundary convention, and $b$
determines $c$. Hence
\begin{equation}\label{dir:eq:hbound}
 h(-q)\le\sum_{1\le a\le\sqrt{q/3}}a<q/3<(q-1)/2 .
\end{equation}
If $r=4$ or $6$, then $E$ must be the quadratic cyclotomic field $\Q(\mu_r)$, since
adjoining any other quadratic field to $E$ would give an even subdegree in $L/E$. Both
choices have class number one, and in class number one the values on principal ideals
already give $L=E(\mu_r)$, a contradiction with $[L:E]=\ell$. Finally $r=2$ also
precludes $h_E=1$, since then $L=E$. Thus $r=2$ and $h_E=\ell$. Genus theory now gives
$D=-q$ for a single odd prime $q\equiv3\pmod4$, and \eqref{dir:eq:hbound} excludes
$q=\ell$.

\emph{Construction and uniqueness.} Choose a nonprincipal prime ideal $P$ of degree one
in $E$ away from $q\ell$, of norm $p$, and write $P^\ell=(\alpha)$; such primes exist
by the prime ideal theorem in ideal classes. Take as $\eta_0$ the Legendre character
on $(\OO_E/(\sqrt{-q}))^\times=\F_q^\times$. Its rational restriction is $\chi_{-q}$,
and its sign at $-1$ ensures that $\psi((a))=\eta_0(a)a$ is well defined. Extend to the
ideal group by choosing $\psi(P)^\ell=\eta_0(\alpha)\alpha$, using that $P$ generates
the cyclic class group. Changing the sign of $\psi(P)$ if necessary, write
$\beta^\ell=\alpha$. The polynomial $X^\ell-\alpha$ is irreducible over $E$: its
reducibility for prime $\ell$ would make $\alpha$ an $\ell$th power by the binomial
irreducibility criterion, forcing $P$ to be principal. The extension of $\psi$
therefore has value field $E(\beta)$ of degree $\ell$ over $E$. Conversely, for any
admissible $\psi$ with finite character of order two, use a representative $xP$ prime
to its conductor. The value $\psi(xP)/x$ has $\ell$th power $\pm\alpha$ and generates
$L$ over $E$. Thus all resulting value fields are conjugate. Since $\beta\bar\beta=p$,
the real element $t=\beta+p/\beta$ is a root of $D_\ell(X,p)-\Tr_{E/\Q}\alpha$. If $t$
were rational, $\beta$ would be at most quadratic over $\Q$ and
$\beta^\ell\in E\setminus\Q$ would force $\beta\in E$. Hence $t$ generates the real
coefficient field of prime degree $\ell$.

\emph{Injectivity.} The prime $q$ can be read off from $K$: it is the unique ramified
rational prime of $K$ away from $\ell$. Indeed $M=E(\zeta_\ell,\beta)$ is Galois over
$\Q$, since it contains the radicals above both $\alpha$ and $\bar\alpha$. Away from
$q\ell$, its radical extension over $E(\zeta_\ell)$ is unramified: the valuation of
$\alpha$ is a multiple of $\ell$ everywhere, and after removing a local $\ell$th power
one has a unit whose $\ell$th-root polynomial has unit discriminant. At $q\ne\ell$,
inertia maps onto the nontrivial inertia of $E/\Q$ and fixes $\zeta_\ell$. An inertia
lift of conjugation sends $\beta$ to $\zeta_\ell^cp/\beta$ for some $c$. On the $\ell$
distinct conjugates
\[
 t_j=\zeta_\ell^j\beta+\zeta_\ell^{-j}p/\beta\qquad(j\bmod\ell)
\]
it acts by $j\mapsto-j-c$, a nonidentity permutation. Hence $q$ ramifies in the
splitting field and in $K$, because the normal closure of a field unramified at $q$ is
unramified there.

The counts are those of the imaginary quadratic fields of class number $7$, $11$ and
$13$ in the unconditional enumeration \cite{Watkins2004Class}.
\end{proof}




\subsection{The discriminant and conductor of the reflected fields}
\label{dir:sec:reflected}
\sources{grh-exponent-six-and-transfers; grh-lean-transfers; grh_transfers}

The fields of Theorems~\ref{dir:thm:cmquintic} and~\ref{dir:thm:cmoddprime} belong to the reflection
families of Cohen, Rubinstein-Salzedo and Thorne
\cite{CohenRubinsteinSalzedoThorne2015Identities,CohenThorneDihedral}, whose reflection
theorem leaves two possible exponents of $\ell$ in the discriminant. This subsection
decides which one occurs, gives the conductor in the Galois closure, and separates the
fields over the two quadratic bases of class number $25$.

Let $D<0$ be a fundamental discriminant and $\ell$ an odd prime coprime to $D$. Put
$E=\Q(\sqrt D)$, $B=E(\zeta_\ell)$ and $r=\dim_{\F_\ell}\Cl(E)[\ell]$. Under these
hypotheses every unit of $E$ is an $\ell$th power.

\begin{proposition}[The reflection family; \stU]\label{dir:prop:reflection}
Let $[\mathfrak a]$ be a class of order $\ell$ in $\Cl(E)$, write
$\mathfrak a^\ell=(\alpha)$, $n=\Nm\mathfrak a$ and $\beta^\ell=\alpha$, and put
$M=B(\beta)$. Let $c$ denote complex conjugation.
\begin{enumerate}[label=\textup{(\roman*)}]
\item The $\ell$ numbers $\zeta_\ell^j\beta+n/(\zeta_\ell^j\beta)$ are real, distinct
 and conjugate. Their splitting field is $M^{\langle c\rangle}$, its Galois group is
 the affine group $\F_\ell\rtimes\F_\ell^\times$, and its cyclic subfield of degree
 $\ell-1$ is $B^{\langle c\rangle}$.
\item The isomorphism class of the field $\Q(\beta+n/\beta)$ depends only on the line
 spanned by $[\mathfrak a]$ in $\Cl(E)[\ell]$, and different lines give different
 normal closures. There are $(\ell^r-1)/(\ell-1)$ such fields.
\end{enumerate}
\end{proposition}

\begin{proof}
We have $\alpha\bar\alpha=n^\ell$. Changing the ideal representative changes $\alpha$
by an $\ell$th power, up to a unit, and every unit is an $\ell$th power. Thus the
group of virtual units
\[
 \set{v\in E^\times:(v)\text{ is an }\ell\text{th power of an ideal}}/E^{\times\ell}
 \simeq\Cl(E)[\ell]
\]
identifies ideal-class lines with radical lines. A radical line remains nonzero and
distinct from the other lines in $B^\times/B^{\times\ell}$, because restriction
followed by the norm raises classes to the power $[B:E]=\ell-1$.

A generator $s$ of $\Gal(M/B)$ sends $\beta$ to $\zeta_\ell\beta$. For a primitive root
$g$ modulo $\ell$ choose $u$ fixing $E$ and $\beta$ and sending $\zeta_\ell$ to
$\zeta_\ell^g$. Complex conjugation $c$ sends $\beta$ to $n/\beta$ and $\zeta_\ell$ to
$\zeta_\ell^{-1}$. These automorphisms generate the Galois group of $M/\Q$, and direct
substitution shows that $c$ is central and $usu^{-1}=s^g$. The $\ell$ real conjugates
of $\beta+n/\beta$ are $\zeta_\ell^j\beta+n/(\zeta_\ell^j\beta)$. They are distinct:
equality for different $j$ would make $\beta^2/n$ an $\ell$th root of unity and hence
$\alpha^2=n^\ell$, forcing $[\alpha]=0$ in $E^\times/E^{\times\ell}$. Their splitting
field is $M^{\langle c\rangle}$, since the affine action of $s$ and $u$ on the indices
is faithful. Its cyclic subfield $B^{\langle c\rangle}$ of degree $\ell-1$ is the
mirror field in the sense of \cite{CohenRubinsteinSalzedoThorne2015Identities}, and
the conjugation rule is their star condition.

Two different ideal-class lines give different normal closures: on adjoining $E$ each
normal closure recovers its Kummer extension $M/B$, and Kummer theory distinguishes the
lines. Different generators of one line give the same normal closure. Its cyclic
complements of order $\ell-1$ are conjugate in the affine group, so they give a single
isomorphism class of fields of degree $\ell$.
\end{proof}

For the following calculation assume that $D=-q$ is an odd prime discriminant and
$q\ne\ell$. With the notation of Proposition~\ref{dir:prop:reflection} set
\[
 K=\Q(\beta+n/\beta),\qquad C=M^{\langle c\rangle},\qquad F=B^{\langle c\rangle},
\]
so that $C$ is the normal closure of $K$ and $C/F$ is cyclic of degree $\ell$. Let
$v\mid\ell$ be a place of $E$, $U=E_v$ and $f=[U:\Q_\ell]$, and let $u$ be the unit
obtained from $\alpha$ by dividing by $\ell^{v(\alpha)}$, with $v(\ell)=1$. Since
$\ell\mid v(\alpha)$, this preserves the local radical class.

\begin{theorem}[The local test at $\ell$; \stU]\label{dir:thm:localtest}
With these assumptions, the discriminant of $K$ and the relative Artin conductor of
any nontrivial character of $C/F$ are given by
\[
 \begin{array}{c|c|c}
 \text{test in }\OO_U/\ell^2 & \abs{\disc(K)} & \text{conductor of }C/F\\ \hline
 u^{\ell^f-1}=1 & q^{(\ell-1)/2}\,\ell^{\ell-2} & (1)\\
 u^{\ell^f-1}\ne1 & q^{(\ell-1)/2}\,\ell^{\ell} & \lambda^2
 \end{array}
\]
where $\lambda$ is the unique prime of $F$ over $\ell$. If $\ell$ splits in $E$, the
two choices of $v$ give the same answer, since $\alpha\bar\alpha=n^\ell$.
\end{theorem}

\begin{proof}
\emph{The local algebra at $\ell$.} We have $KE=E(\beta)$, by $[E(\beta):E]=\ell$ and
$\beta^2-(\beta+n/\beta)\beta+n=0$. Extension of scalars gives
\[
 (K\otimes_\Q\Q_\ell)\otimes_{\Q_\ell}U\simeq U[X]/(X^\ell-u).
\]
The unramified extension $U/\Q_\ell$ commutes with the formation of the maximal order
of an \'etale algebra: on each local factor its base change is a product of rings of
integers of unramified extensions. Hence the discriminant valuation can be read after
this scalar extension.

Decompose a unit of $U$ into a Teichm\"uller root of unity of order dividing
$\ell^f-1$ and a principal unit. The $\ell$th power map is an automorphism of the
former group. On the latter the logarithm identifies $1+\ell\OO_U$ with $\ell\OO_U$,
and taking $\ell$th powers has image $1+\ell^2\OO_U$. Thus $u$ is an $\ell$th power in
$U$ exactly when $u^{\ell^f-1}=1$ in $\OO_U/\ell^2$. In that case the algebra is
$U\times U(\zeta_\ell)$, with a tame, totally ramified factor of degree $\ell-1$, and
its discriminant exponent is $\ell-2$.

In the other case take $z$ to be any lift of the unique residue root of
$X^\ell-\bar u$. Then $v(z^\ell-u)=1$, since a solution modulo $\ell^2$ would pass the
power test. The shifted polynomial $(X+z)^\ell-u$ is Eisenstein. Its root is a
uniformizer and generates the ring of integers; the derivative of the original
polynomial at its root is $\ell$ times a unit, and the norm of this value has valuation
$\ell$. This gives the wild exponent $\ell$.

\emph{The remaining primes.} At a finite prime of $B$ away from $\ell$ the valuation of
$\alpha$ is a multiple of $\ell$. After removing a local $\ell$th power one adjoins a
root of a polynomial $X^\ell-v_0$ with $v_0$ a unit and unit discriminant, hence an
unramified extension. Consequently ramification away from $\ell$ can occur only at $q$.
There inertia has order two in $M/\Q$, fixes $\zeta_\ell$, and acts on the $\ell$ real
roots by $j\mapsto h-j$ for some $h$. It has one fixed point and $(\ell-1)/2$
transpositions, so tame inertia contributes $(\ell-1)/2$ to the discriminant exponent
of $K$ at $q$. This inertia also injects into $\Gal(F/\Q)$, so $C/F$ is unramified
there.

\emph{The conductor.} At $\ell$, inertia acts on $F$ through the full cyclotomic group,
since $E$ is unramified there and the cyclotomic group acts faithfully on $F$. Thus
$F/\Q$ is totally and tamely ramified at $\ell$. In the power case, adjoining $U$ to
the local splitting field leaves it inside $U(\zeta_\ell)$, of absolute degree at most
$2(\ell-1)$. Its local degree over $F_\lambda$ divides $\ell$, so $\lambda$ splits
completely in $C/F$.

In the wild case, over $U$ the splitting field contains the totally ramified radical
extension of degree $\ell$ and the totally ramified cyclotomic extension of degree
$\ell-1$; indeed the normal closure of $K$ after this base change is the splitting
field of the isomorphic radical algebra. The ramification index is divisible by both
degrees, and its upper bound $\ell(\ell-1)$ shows that the entire affine group is
inertia already over $\Q_\ell$. Use $K$ locally in the discriminant tower: the
extension $C/K$ is tame of degree $\ell-1$, so
\[
 v_\ell(d_{C/\Q})=(\ell-1)\ell+(\ell-2)=\ell^2-2 .
\]
Using $F$ instead gives $v_\lambda(d_{C/F})=\ell^2-2-\ell(\ell-2)=2(\ell-1)$, as all
residue degrees are one. The conductor--discriminant formula for the cyclic group of
prime order gives conductor exponent two for each of its $\ell-1$ nontrivial
characters.
\end{proof}

\begin{corollary}[The quintic CM coefficient fields; \stU]\label{dir:cor:quinticdisc}
For the twenty-five fields of \cref{dir:tab:quintic} the discriminant is
$q^2\,5^{v_5}$ with $v_5$ as tabulated. Exactly six of them are tame at five, those
with $q=347,443,739,1051,1123,1723$. The field $\Q(\zeta_{11})^+$ has discriminant
$11^4$. The discriminant of the polynomial $P_{p,a}$ is $5^5(4p^5-a^2)^2=5^5q^2b^4$,
so the index of $\Z[X]/(P_{p,a})$ in the maximal order is $5b^2$ in the tame case and
$b^2$ in the wild case.
\end{corollary}

\begin{proof}
For each row, $p$ is a split rational prime, $4p^5=a^2+qb^2$ and $p\nmid a$; with
$\alpha=(a+b\sqrt{-q})/2$ the ideal $(\alpha)$ is the fifth power of a prime above $p$,
and as $p<q/4$ a generator of norm $p$ cannot exist. The polynomial $P_{p,a}$ follows
by expanding $(\beta+p/\beta)^5$ and using $\beta^5+\bar\beta^5=a$. Apply
Theorem~\ref{dir:thm:localtest} with $\ell=5$. The checker enumerates all $625$ potential
roots modulo $25$ in the quadratic ring. When $p=5$ the quadratic algebra at
five splits; one component of $\alpha$ is a unit and the other has valuation five, so
modulo $25$ the latter is zero, the unit component decides the test, and its
reciprocal times $5^5$ gives the other local class. The discriminant of $P_{p,a}$ is
obtained from a resultant with $P'_{p,a}$, or from the critical values $\pm2p^{5/2}$
of the Dickson polynomial. For $\Q(\zeta_{11})^+$ use the conductor--discriminant
formula for its four nontrivial characters. A separate computation of maximal orders
by the Round Two algorithm verifies every field discriminant and order index.
\end{proof}
\dirfiles{\artifact{grh-exponent-six-and-transfers/checks/reflection.py},
\artifact{grh-exponent-six-and-transfers/checks/maximal_orders.py},
\artifact{grh-exponent-six-and-transfers/data/reflected_fields.csv}}

For example, for $q=347$ one may take $p=3$, $a=25$, $b=1$. Modulo $25$ we have
$(2\sqrt{-347})^5=(25+\sqrt{-347})/2$, since $32\cdot347^2\equiv13\pmod{25}$; the field
discriminant is $5^3\cdot347^2=15\,051\,125$ and the index is five. For $q=47$, $p=2$,
$a=9$, $b=1$, the field discriminant is $5^5\cdot47^2=6\,903\,125$ and $P_{2,9}$
generates the full ring of integers.

\begin{proposition}[Two class groups of rank two; \stU]\label{dir:prop:ranktwo}
Over each of $D=-12451$ and $D=-37363$, whose class groups are $C_5^2$, the six lines
of $\F_5^2$ give six pairwise nonisomorphic reflected quintic fields, with defining
polynomials $P_{p,a}$ as in \cref{dir:tab:ranktwo}. Exactly one line over each base is
tame at five, and the five wild fields over one base have the same discriminant.
\end{proposition}

\begin{proof}
A projective label certifies that its row represents a distinct line. For each base
choose two triples $(s,r,z)$ in which $s\equiv1\pmod5$ is prime, $r^2\equiv D\pmod s$
and $z$ has order five modulo $s$. Reducing $\alpha$ by $\sqrt D\mapsto r$ and raising
to the power $(s-1)/5$ gives $z^e$ with $e\in\F_5$; the two exponents form the label
$[e_1:e_2]$. All generators in the table are units at the probe primes. A relation
between their classes modulo fifth powers would remain a relation between the labels,
because the residual fifth-power characters are multiplicative. The six
nonproportional labels prove exhaustion by the known rank. The triples are
$(181,20,59)$, $(211,93,107)$ for the first base and $(151,38,59)$, $(241,76,87)$ for
the second. The values of $v_5$ come from Theorem~\ref{dir:thm:localtest}, and the maximal
order computation cited above verifies all twelve rows.
\end{proof}

\begin{table}[ht]
\centering\small
\begin{tabular}{rrrrr@{\qquad}rrrrr}
\toprule
\multicolumn{5}{c}{$D=-12451$} & \multicolumn{5}{c}{$D=-37363$}\\
$p$ & $a$ & $b$ & label & $v_5$ & $p$ & $a$ & $b$ & label & $v_5$\\
\midrule
5 & 7 & 1 & $[2\!:\!0]$ & 3 & 11 & 779 & 1 & $[3\!:\!1]$ & 5\\
7 & 132 & 2 & $[1\!:\!1]$ & 5 & 13 & 942 & 4 & $[1\!:\!3]$ & 5\\
11 & 577 & 5 & $[2\!:\!3]$ & 5 & 23 & 4607 & 11 & $[3\!:\!3]$ & 5\\
13 & 1134 & 4 & $[1\!:\!3]$ & 5 & 37 & 13131 & 53 & $[4\!:\!1]$ & 5\\
53 & 23339 & 301 & $[0\!:\!1]$ & 5 & 43 & 7677 & 119 & $[0\!:\!1]$ & 3\\
109 & 79839 & 2105 & $[2\!:\!4]$ & 5 & 79 & 96628 & 282 & $[3\!:\!0]$ & 5\\
\bottomrule
\end{tabular}
\caption{The twelve reflected quintic fields over the two imaginary quadratic fields
with class group $C_5^2$: $4p^5=a^2+\abs Db^2$, defining polynomial $P_{p,a}$.}
\label{dir:tab:ranktwo}
\end{table}

\begin{corollary}[The reflected quintic fields of exponent five; \stG]
\label{dir:cor:reflected37}
There are exactly $27$ imaginary quadratic fields of class-group exponent five, and
their reflection families for $\ell=5$ contain exactly $37$ quintic fields: the $25$
fields of \cref{dir:tab:quintic} and the $12$ fields of \cref{dir:tab:ranktwo}. Each of
them is a totally real field of discriminant $\abs D^2\,5^3$ or $\abs D^2\,5^5$, and
these are all the starred fields of \cite{CohenRubinsteinSalzedoThorne2015Identities}
with those discriminants.
\end{corollary}

\begin{proof}
Elsenhans, Kl\"uners and Nicolae
\cite[Theorem~1 and Table~1]{Elsenhans2020Imaginary} find $27$ imaginary
quadratic fields of class-group exponent five up to $3.1\cdot10^{20}$, and their
Lemma~1 and Theorem~3 bound every absolute discriminant of exponent five by
$2.3\cdot10^{20}$ under the quadratic least-split-prime estimate (their Table~2). The
required statement for $\zeta_E(s)=\zeta(s)L(s,\chi_D)$ follows from
\cref{src:thm:dgrh}. Their data give $25$ groups $C_5$ and two groups $C_5^2$, at
$D=-12451$ and $D=-37363$, so the count is $25+2(5^2-1)/(5-1)=37$ by
Proposition~\ref{dir:prop:reflection}. By Theorem~\ref{dir:thm:localtest} the fields lie among the
totally real starred fields of discriminant $\abs D^{(\ell-1)/2}\ell^{\ell-2}$ or
$\abs D^{(\ell-1)/2}\ell^{\ell}$, and Theorem~1.2 and Lemma~2.8 of
\cite{CohenRubinsteinSalzedoThorne2015Identities} give exactly
$(\ell^r-1)/(\ell-1)$ starred fields in the union of these two sets. The construction
therefore exhausts the family; the tables assign a specific alternative and a defining
polynomial to each ideal-class line.
\end{proof}

The weight-two CM classification of Theorem~\ref{dir:thm:cmquintic} uses only the $25$ cyclic
class groups and the cyclotomic exception. The twelve fields of
\cref{dir:tab:ranktwo} are the subject of a joint splitting test in \cref{sec:grp}.

\subsection{A uniform isogeny bound for two families of sextic fields}
\label{dir:sec:isogeny}
\sources{grh-arithmetic-transfers; dirichlet-grh-research}

Let $E$ be an elliptic curve over a number field $K$. A cyclic isogeny of prime degree
$p$ defined over $K$ gives a Galois-stable line in $E[p]$; write $\lambda$ for the
character describing the action on the line and $\chi_p$ for the mod-$p$ cyclotomic
character. We use \emph{Momose type 2} in the sense of Banwait and Derickx, with
$\lambda^{12}=\chi_p^6$ and $p\equiv3\pmod4$ as in their type~2 sieve
\cite{BanwaitDerickx2025Explicit}. A bound for such $p$ that is independent of the
discriminant of $K$ is the uniformity question; Banwait and Derickx treat cubic fields
and ask how their method might treat fields of even degree
\cite{BanwaitDerickx2023Towards}. The arithmetic of this subsection rests on a statement
about residue degrees that is also used in \cref{dir:sec:heavenly}.

\begin{lemma}[Odd residue degrees; \stU]\label{dir:lem:oddresidue}
Let $K$ be a number field with normal closure $N/\Q$, and let $G=\Gal(N/\Q)$ act on
the set $X=G/H$ of embeddings of $K$, where $H=\Gal(N/K)$. If every element of $G$ has
an odd cycle of length at most $r$ on $X$, then every rational prime $p$ has a place
$v$ of $K$ above it of odd residue degree $f_v\le r$. This includes the primes
ramified in $N$.
\end{lemma}

\begin{proof}
Choose a place of $N$ above $p$, its decomposition group $D$, its inertia group
$I\triangleleft D$, and a lift $u\in D$ of the Frobenius generator of $D/I$. The places
of $K$ above $p$ correspond to the $D$-orbits on $X$. In the orbit of $x$, with
stabilizer $D_x$, the inertia orbits are permuted transitively by $u$, and their number
is $[D:ID_x]=f_v$: the local degree $[D:D_x]$ of the subfield of $N_w$ fixed by $D_x$,
divided by its ramification index $[I:I\cap D_x]$, gives this index. Thus the length of
the $u$-cycle through $x$ is a multiple of $f_v$. Choose $x$ on an odd cycle of $u$ of
length at most $r$.
\end{proof}

\begin{lemma}[The two sextic actions; \stU]\label{dir:lem:sextic}
Consider $A_4$ acting on the six edges of a tetrahedron and $A_5$ acting on its six
Sylow $5$-subgroups. The numbers of elements of each cycle shape are
\begin{center}
\begin{tabular}{lrrrr}
 & $1^6$ & $1^22^2$ & $3^2$ & $1\,5$\\
\hline
$A_4$ on edges & 1 & 3 & 8 & 0\\
$A_5$ on Sylow $5$-subgroups & 1 & 15 & 20 & 24
\end{tabular}
\end{center}
In particular every element has a cycle of length $1$ or $3$. Each of the two groups
has a unique transitive action on six points up to equivalence, the $A_5$ action is
primitive, and among the sixteen transitive permutation groups of degree six these two
are exactly the ones in which every element has an odd cycle.
\end{lemma}

\begin{proof}
The three double transpositions of $A_4$ each fix two edges of the tetrahedron, and an
element of order three has two orbits of length three on the edges. For $A_5$, each of
the six point stabilizers has order ten and contains five involutions; the fifteen
involutions are conjugate, so each fixes $6\cdot5/15=2$ points. An element of order
five fixes its unique Sylow subgroup, and an element of order three fixes none. For
$A_4$ the stabilizer has order two and all such subgroups are conjugate. In $A_5$ a
subgroup of order ten has a unique Sylow $5$-subgroup, hence is its normalizer, and all
these normalizers are conjugate; the stabilizer is maximal, since the action is doubly
transitive (the $5$-cycle in the stabilizer permutes the other five points). The
statement about the sixteen transitive groups is an exhaustive computation.
\end{proof}
\dirfiles{\artifact{dirichlet-grh-research/heavenly/check_heavenly.py},
\artifact{dirichlet-grh-research/heavenly/check_heavenly.txt}}

The LMFDB number fields \code{6.6.284124736.1} and \code{6.2.27071209.2} have these two
actions.

The necessary condition used below is Proposition~6.1 of
\cite{BanwaitDerickx2025Explicit}: if $p$ is a Momose type~2 isogeny prime for $K$,
then for every rational prime $q$ that splits in $\Q(\sqrt{-p})$ and every prime of $K$
above $q$ of odd residue degree $f$,
\begin{equation}\label{dir:eq:cc}
 4q^f\ge p\qquad\text{or}\qquad p\mid q^{2f}+q^f+1 .
\end{equation}
The second alternative matters for potentially multiplicative reduction. There is no
assumption that $q$ be unramified in $K$.

\begin{theorem}[The strict local sieve; \stG]\label{dir:thm:strictcc}
For a prime $p\equiv3\pmod4$, call a rational prime $w$ an \emph{excluding witness} if
\[
 4w^3<p,\qquad\legendre wp=1,\qquad p\nmid(w^2+w+1)(w^6+w^3+1).
\]
The primes $p\equiv3\pmod4$ with no excluding witness are exactly the following $66$:
\[
\begin{gathered}
3,7,11,19,23,31,43,59,67,83,107,139,163,211,283,307,331,379,499,523,\\
547,643,787,883,907,1123,1723,1747,2083,2347,2683,3067,3187,3907,\\
4027,4243,4363,4603,5107,5443,5923,6043,6427,6763,7723,8563,8803,\\
9067,10627,10987,13267,14107,14347,14683,15667,16963,22003,26203,\\
32323,34483,36523,37123,89083,90787,222643,253507 .
\end{gathered}
\]
\end{theorem}

By quadratic reciprocity, $\legendre wp=1$ is equivalent to $w$ splitting in
$\Q(\sqrt{-p})$, also for $w=2$.

\begin{proof}
Trial of the eligible auxiliary primes gives the list up to $253507$.

\emph{The range $253507<p\le11\cdot10^{12}$.} An independent sieve with the strict
integer inequality finds, or certifies the existence of, an excluding witness for every prime in this range. It
checks $264911$ primes directly up to $4\cdot127^3=8193532$. Above this threshold all
primes $q\le127$ have admissible size. The program separately checks the eight primes
$p\equiv3\pmod4$ in its upper range that divide one of the two polynomials for some
$q\le127$. For all remaining primes, a square among these auxiliaries is a witness. If
every auxiliary is a nonsquare, the conditions at $2,3,5,7,11,13,17$ restrict $p$ to
$1440$ residue classes modulo $2042040$. For every odd $q$, quadratic reciprocity gives
$\legendre qp=\legendre{-p}q$; in each class the further tests for $19\le q\le127$ are
therefore periodic masks on an arithmetic progression, a mask discarding the multiples
of $q$ and the integers with $\legendre{-p}q=1$. Exactly $327$ integers in the upper
range survive all the masks. Trial division proves that $110$ of them are prime, and
all $110$ have an excluding witness.

\emph{The range $p>T=10^{11}$.} Put $y=\log p$ and $M(y)=1.881y+6.18$. For
$p>T>e^{25}$, Theorem~3 of \cite{Elsenhans2020Imaginary}, which rests on
Bach and Sorenson \cite{Bach1996Explicit} and on the Riemann hypothesis for the
two Dirichlet factors of $\zeta_{\Q(\sqrt{-p})}(s)=\zeta(s)L(s,\chi_{-p})$, gives a
prime $q\le M(y)^2$ that splits in $\Q(\sqrt{-p})$; the hypothesis is supplied by
\cref{src:thm:dgrh}. Here $4q^3<p$: indeed $M(\log T)<54$ and $4\cdot54^6<T$, and
$M(y)^6e^{-y}$ decreases for $y\ge\log T$ by logarithmic differentiation. If $y\ge58$,
the same argument applied to the exponents $12$ and $6$, using $M(58)<116$ and
$116^{12}+116^6+1<e^{58}$, yields $q^6+q^3+1<p$, and $q$ is already an excluding
witness.

In the remaining case $q\le116^2$ and $q^2+q+1<T$. The only obstruction to using $q$ is
$p\mid N_q=q^6+q^3+1$. For each of the $1594$ primes $q\le116^2$ the supplied
certificate decomposes $N_q=\prod_id_i^{e_i}$. A prime $p>T$ dividing $N_q$ divides some $d_i>T$. For every
such entry the certificate supplies a prime $w\le43$ and an integer $z$ such that
\[
 z^2\equiv w\pmod{d_i},\qquad
 \gcd\bigl(d_i,\,w(w^2+w+1)(w^6+w^3+1)\bigr)=1 .
\]
Reducing modulo $p$ gives a nonzero square $w$ and both required nondivisibilities, and
since $4w^3<T$ this $w$ is an excluding witness. The argument works even when $d_i$ is
composite, and no primality claim about the large divisors enters. The verifier checks
every product, square, gcd and endpoint inequality with integer or rational
arithmetic; it covers $1126$ distinct large divisors. The two overlapping ranges prove
the classification.
\end{proof}
\dirfiles{\artifact{grh-arithmetic-transfers/checks/isogeny_sieve.py} (the list up to
$253507$), \artifact{grh-arithmetic-transfers/checks/cubic_tail.cpp} (the middle range),
\artifact{grh-arithmetic-transfers/checks/large_certificate.json} and
\artifact{grh-arithmetic-transfers/checks/verify_large.py} (the certificate for large
primes and its verifier)}

\begin{remark}\label{dir:rem:strict}
The strict inequality $4q^3<p$ must be retained even when $p$ lies just below $4q^3$.
The loop in Algorithm~A.8 of \cite{BanwaitDerickx2023Towards}, and the authors'
PARI/GP implementation, also admit $q\le1+(p/4)^{1/3}$, which can include an extra
prime. For example $p=89083$ has no excluding witness; the extra prime $q=29$ has
$\legendre{29}{p}=1$ and neither polynomial vanishes modulo $p$, yet
$4\cdot29^3=97556>p$. The endpoint $253507$ agrees with the one reported in
\cite[Appendix~A, Example~A.13]{BanwaitDerickx2023Towards} for the test with residue
degrees $1,3$. A second run of the tail program, with its upper endpoint raised to
$3\cdot10^{13}$, leaves $782$ integers after the masks, of which $246$ are prime, each
with an excluding witness (\artifact{dirichlet-grh-research/heavenly/cubic_tail.txt}).
The prime $p=3$ has no $q$ with $4q<p$ and passes vacuously.
\end{remark}

\begin{theorem}[Sextic fields; \stG]\label{dir:thm:sexticisogeny}
Let $K$ be a sextic field whose normal closure acts on the six embeddings of $K$ as
$A_4$ on the edges of a tetrahedron or as $A_5$ on its Sylow $5$-subgroups. Then every
odd Momose type~2 isogeny prime for $K$ belongs to the set of $66$ primes of
Theorem~\ref{dir:thm:strictcc}, whose largest member is $253507$.
\end{theorem}

\begin{proof}
Let $p\equiv3\pmod4$ be a prime outside the list, and let $w$ be an excluding witness.
Then $w$ splits in $\Q(\sqrt{-p})$. By Lemmas~\ref{dir:lem:oddresidue} and~\ref{dir:lem:sextic} some
prime of $K$ above $w$ has residue degree $f\in\set{1,3}$. Since $4w^f\le4w^3<p$ and
$p$ divides neither $w^2+w+1$ nor $w^6+w^3+1$, both alternatives of \eqref{dir:eq:cc}
fail, so $p$ is not a Momose type~2 isogeny prime for $K$.
\end{proof}

The analytic input concerns $\Q(\sqrt{-p})$ only, which is why Dirichlet GRH suffices
although the normal closure is nonabelian, and for $A_5$ nonsolvable.

\begin{proposition}[A general uniform bound; \stG]\label{dir:prop:isogenygeneral}
Suppose that every element of a transitive permutation group $G$ acting on $G/H$ has
an odd cycle of length at most $m$, and let $K$ range over the number fields whose
normal closure acts in this way on the embeddings of $K$. Put $U(p)=(4\log p+10)^2$.
Every Momose type~2 isogeny prime $p$ for such a field satisfies
\[
 p\le\max\set{4U(p)^m,\ U(p)^{2m}+U(p)^m+1},
\]
an effective bound depending only on $m$. More generally, let $F/\Q$ be a fixed
abelian extension contained in the normal closure $N$, and suppose only that every
element of $\Gal(N/F)$ has an odd cycle of length at most $m$. Then the Momose type~2
isogeny primes of these fields are bounded in terms of $m$ and $F$.
\end{proposition}

\begin{proof}
The explicit bound of Bach and Sorenson \cite{Bach1996Explicit}, as used in
\cite[Appendix~A]{BanwaitDerickx2023Towards}, gives a prime $q\le U(p)$ that splits in
$\Q(\sqrt{-p})$; its hypothesis is the Riemann hypothesis for $\zeta(s)L(s,\chi_{-p})$.
By Lemma~\ref{dir:lem:oddresidue} some prime of $K$ above $q$ has odd residue degree
$f\le m$. If both displayed expressions were smaller than $p$, both alternatives in
\eqref{dir:eq:cc} would fail for this prime. The right side is a fixed polynomial in
$\log p$.

For the variant, apply the least-split-prime estimate to the abelian field
$J=F\,\Q(\sqrt{-p})$, whose zeta function is a product of Dirichlet $L$-functions by
Proposition~\ref{dir:prop:abelian}. Its degree is at most $2[F:\Q]$ and its absolute discriminant
is at most $\abs{D_F}^2p^{[F:\Q]}$, by the discriminant inequality for a compositum.
The resulting prime $q\ll_F(1+\log p)^2$ splits in both fields. Its decomposition
group in $N$ acts trivially on $F$, so the argument of Lemma~\ref{dir:lem:oddresidue}
applies, even if $q$ ramifies in $N/F$.
\end{proof}

For example, the action of $S_4$ on the edges of the tetrahedron has the property of
the variant with $m=3$ once its quadratic sign field is fixed, which gives a bound
uniform over these sextic fields with fixed sign field.

The set of $66$ primes cannot be shortened by the local condition alone.

\begin{theorem}[Sharpness of the local obstruction; \stG]\label{dir:thm:isogenysharp}
There is a sextic field $K$ with the tetrahedral $A_4$ action, and there is a sextic
field with the primitive $A_5$ action, with the following property: among the primes
$p\equiv3\pmod4$, the condition \eqref{dir:eq:cc} holds for every rational prime $q$
that splits in $\Q(\sqrt{-p})$ and every prime of the field above $q$ of odd residue
degree exactly when $p$ is one of the $66$ primes of Theorem~\ref{dir:thm:strictcc}. That
\eqref{dir:eq:cc} holds for these $66$ primes is unconditional.
\end{theorem}

\begin{proof}
Let $S$ be the union, over the $66$ candidates $p$, of the sets of primes $q<p/4$ that
split in $\Q(\sqrt{-p})$. It comprises $5093$ primes.

\emph{A cyclic cubic field in which every prime of $S$ is inert.} Consider
\[
 u_t(X)=X^3-tX^2-(t+3)X-1,\qquad\disc(u_t)=(t^2+3t+9)^2 .
\]
For every prime $q$, the values $0$ and $-1$ are never roots modulo $q$. Each of the
other $q-2$ residue classes $x$ gives at most one $t$ with $u_t(x)=0$, since the
coefficient of $t$ is $-x(x+1)$. Thus there is some $t$ modulo $q$ with $u_t$
irreducible modulo $q$. Choose these parameters simultaneously for all $q\in S$ by the
Chinese remainder theorem. Then $u_t$ is irreducible over $\Q$, its splitting field
$C=\Q(\alpha)$ is cyclic cubic, and every $q\in S$ is unramified and inert, because the
reduction is separable and irreducible. One can take $\sigma(\alpha)=-1/(\alpha+1)$ as
a generator of $\Gal(C/\Q)$, as substitution into $u_t$ verifies.

\emph{The tetrahedral extension.} Choose a rational prime splitting completely in $C$
and use weak approximation to find $\beta\in C^\times$ with valuations $(1,0,0)$ at its
three primes. Set $a=\beta/\sigma\beta$. The classes of $a$ and $\sigma a$ are
independent in $C^\times/(C^\times)^2$, since their valuation vectors modulo two are
two distinct nonzero vectors of even weight on the three primes. The relation
$a\,\sigma a\,\sigma^2a=1$ shows that $N=C(\sqrt a,\sqrt{\sigma a})$ is Galois over
$\Q$. Its quotient of order three permutes the three nontrivial square classes
cyclically, and a subgroup of order three exists by Cauchy's theorem, so the group is
the nontrivial semidirect product $V_4\rtimes C_3\cong A_4$. The sextic subfield
$K=C(\sqrt a)$ has the tetrahedral edge action.

\emph{Verification of \eqref{dir:eq:cc}.} Fix one of the $66$ primes $p$ and a rational
prime $q$ that splits in $\Q(\sqrt{-p})$. If $q\ge p/4$, the first alternative holds
for every $f$. Otherwise $q\in S$, and every residue degree of $K$ above $q$ is a
multiple of $3$ and at most $6$, so the only odd degree to consider is $3$. If
$4q^3\ge p$ the first alternative holds. If $4q^3<p$, then $q^2+q+1<p$, and since $p$
has no excluding witness we must have $p\mid q^6+q^3+1$, the second alternative with
$f=3$. This proves \eqref{dir:eq:cc} for $K$ and all $66$ primes. (The finite check
shows that in fact $4q^3\ge p$ for every such pair $(p,q)$.) Conversely, by
Theorem~\ref{dir:thm:sexticisogeny} and its proof, \eqref{dir:eq:cc} fails for every other
prime $p\equiv3\pmod4$; this is the only step that uses Theorem~\ref{dir:thm:strictcc}.

\emph{The primitive family.} We construct an $A_5$-extension in which every $q\in S$ is
unramified with Frobenius a $3$-cycle in the natural quintic action; in the sextic
action its cycle lengths are $3,3$, and the verification above applies unchanged. Let
$X$ parametrize monic quintics $f$ together with $y\ne0$ satisfying $y^2=\disc(f)$. Its
function field is $\Q(x_1,\dots,x_5)^{A_5}$: the elementary symmetric functions of the
roots generate the $S_5$-invariant field and the Vandermonde product generates its
quadratic extension. Maeda proved that this field is purely transcendental
\cite{Maeda1989Noether}. Consequently $X$ has weak approximation at finite sets of
primes: a nonempty Zariski open subset of $X$ is isomorphic to an open subset of affine
$5$-space, every nonempty open subset of $X(\Q_q)$ meets this chart since $X$ is smooth
and geometrically integral, and ordinary weak approximation in the five coordinates
suffices. At every $q\in S$ prescribe squarefree reduction of type $(3,1,1)$; such a
polynomial exists even over $\F_2$, by taking an irreducible cubic and two distinct
linear factors. At two further primes prescribe types $(3,1,1)$ and $(5)$. Each
prescription defines a nonempty open subset of the local points of $X$: a monic
integral lift of the chosen reduction has unramified splitting field with Frobenius of
the indicated cycle type, this permutation is even, so the signed Vandermonde product
of its roots belongs to $\Q_q$ and supplies $y$. Weak approximation gives one rational
quintic realizing all prescriptions. Its group is contained in $A_5$ and has order
divisible by $15$. A subgroup of order $30$ would contradict simplicity, and a subgroup
of order $15$ would have a normal Sylow $5$-subgroup, whereas the normalizer of this
subgroup in $A_5$ has order $10$. Hence the splitting field has group $A_5$.
\end{proof}

\begin{remark}[Explicit witnesses]\label{dir:rem:sexticwitness}
The tetrahedral field of Theorem~\ref{dir:thm:isogenysharp} is explicit. The union $S$ omits
the prime $r=9343$. Take $t\equiv0\pmod{r^2}$ in addition to the inertness conditions.
The polynomial $X^3-3X-1$ has the simple roots $755,8688,9243$ modulo $r$, and its
value at $755$ has $r$-adic valuation one, so $\beta=\alpha-755$ has the valuations
required above. With the two integers $A,B$ of the first parameter file listed below,
\[
 t=A/B,\qquad u_t(\alpha)=0,\qquad
 K=\Q\Bigl(\alpha,\sqrt{(\alpha-755)/(-1/(\alpha+1)-755)}\Bigr).
\]
The largest absolute coefficient of $Bu_t(X)$ has $10797$ decimal digits. The verifier
rebuilds the union and checks the parameter at each of its primes, including $2$.

For the single endpoint $p=253507$ the set of split primes $q<p/4$ has $3152$ members,
from $41$ to $63361$, and already $4\cdot41^3>p$. Here the parameter is obtained as
follows. To shorten the defining polynomial, the admissible residues are first combined
into $t_0$ modulo $M=17^2\prod_qq$, the product over these $3152$ primes, with
$t_0\equiv0\pmod{17^2}$. The Euclidean algorithm gives a short pair $(A,B)$ with
$A\equiv t_0B\pmod M$ and $\gcd(B,M)=1$, and $t=A/B$; since the denominator is a unit at
each prescribed prime, reduction and irreducibility work as before. Modulo $17$, $u_t$
has the three simple roots $3,4,10$ and $u_t(3)\equiv17\pmod{17^2}$, so
$\beta=\alpha-3$ has the required valuations. The two integers are in the second
parameter file. The largest absolute coefficient of $BX^3-AX^2-(A+3B)X-B$ has $6796$ decimal digits; a monic integral
polynomial for the same cubic field is $Y^3-AY^2-B(A+3B)Y-B^3$, with largest
coefficient of $20382$ digits. This cyclic cubic field is also a witness for the local
compatibility question of \cite[Remark~A.14]{BanwaitDerickx2023Towards}, whose cubic
example has coefficients exceeding $27380$ decimal digits and enforces inertness at
every prime below $p/4$; the weaker splitting prescriptions used here suffice for the
same necessary condition.
\dirfiles{\artifact{grh-arithmetic-transfers/checks/simultaneous_parameter.txt},
\artifact{grh-arithmetic-transfers/checks/sextic_parameter.txt} (the two parameter
files), \artifact{grh-arithmetic-transfers/checks/simultaneous_witness.py} (the
verifier)}
\end{remark}

The argument of Proposition~\ref{dir:prop:isogenygeneral} applies to an infinite sequence of
even degrees. Let $N/\Q$ be a Galois extension with group $A_n$, $n\ge4$. Fix a
$k$-element subset of the $n$ letters and let $K$ be the fixed field of its setwise
stabilizer, so that $[K:\Q]=\binom nk$; by complementation we may take
$1\le k\le n/2$. Orbits of a permutation are taken on the set of all $k$-subsets.

\begin{theorem}[Alternating-group subset fields; \stU]\label{dir:thm:subsets}
Every element of $A_n$ has an odd orbit on the $k$-subsets if and only if either
$\binom nk$ is odd, or
\begin{equation}\label{dir:eq:subsetcases}
 n\in\set{2^a,2^a+1},\qquad k=2^{a-1},\qquad a\ge2 .
\end{equation}
In the cases \eqref{dir:eq:subsetcases}, the maximum over $g\in A_n$ of the length of
the shortest odd orbit of $g$ is exactly the largest odd integer $m\le n$. If the
criterion fails, a set of unramified primes of positive density has only even residue
degrees in $K$.
\end{theorem}

\begin{proof}
A permutation $g$ has an odd orbit on subsets exactly when the Sylow $2$-subgroup of
$\langle g\rangle$ fixes a subset; this follows by taking the index of the stabilizer
in the cyclic group. It therefore suffices to test the even permutations of $2$-power
order. Their cycle lengths form partitions of $n$ into powers of two with an even
number of even parts, and the sizes of the fixed subsets are exactly the subset sums
of the parts.

Write $\mathcal B$ for the set of powers of two in the binary expansion of $n$. Every
partition into powers of two refines $\mathcal B$, as one sees by merging equal parts
until all parts are distinct. Splitting a part leaves every previously attainable
subset size attainable. Therefore every submask of $n$, that is, every sum of a subset
of $\mathcal B$, always occurs. These sizes have $\binom nk$ odd, by expanding
$(1+X)^n=\prod_{2^j\in\mathcal B}(1+X^{2^j})$ in $\F_2[X]$.

Suppose $k$ is not a submask. If $\mathcal B$ itself describes an even permutation, it
excludes $k$. Otherwise every even partition is a proper refinement, and replacing any
one part $2^j>1$ of $\mathcal B$ by two halves already gives an even permutation. Thus
$k$ must be a subset sum after \emph{every} possible single split; this is sufficient
too, because a refinement starts with one of those splits. For a given $j$, a new sum
must use exactly one of the two halves. If $2^{j-1}\in\mathcal B$ already, the two
copies used in such a sum either leave the binary bit at $j-1$ or carry into the
vacated bit $j$, and all sums are old. If $2^{j-1}\notin\mathcal B$, a new sum has
exactly one bit outside $\mathcal B$, at $j-1$. Two different choices of $j$ therefore
cannot allow the same new sum. The only remaining $n$ have a single binary part
greater than one, with an optional part $1$. For $n\ge4$ the single split allows
exactly the additional sizes $2^{a-1}$ and, when $n$ is odd, its complement
$2^{a-1}+1$. This proves the equivalence. For an excluded pair $(n,k)$ the Chebotarev
density theorem realizes the conjugacy class of the excluding element as a Frobenius
class with positive density.

For the orbit length put $k=2^{a-1}$ and write each cycle length of $g$ as $2^bu$ with
$u$ odd. The $2$-part $2^b$ divides $k$ in both cases of \eqref{dir:eq:subsetcases}:
a cycle violating this would have length $2^a$ and at most one point outside it,
making the permutation odd. Sort the cycles by decreasing $2^b$. Take whole cycles
until the next cycle would reach or exceed $k$ letters. The remaining number of
letters required is a multiple of the $2^b$ of this cycle. Its $u$th power splits the
cycle into $u$ cycles of length $2^b$; take the required number of these. The
resulting $k$-set is fixed by $g^u$, so its orbit has odd length dividing $u\le m$.

When $n=2k$, a cycle of length $m=2k-1$ with a fixed letter attains $m$: a $k$-set
selects $k$ or $k-1$ letters in the long cycle, both numbers are coprime to $m$, and
its stabilizer in this regular cyclic action has order dividing the number selected,
hence is trivial. When $n=2k+1$ the $n$-cycle attains $m=n$ by the same coprimality
argument. Both witness permutations are even.
\end{proof}

\begin{corollary}[Uniform bounds in arbitrarily large even degree; \stG]
\label{dir:cor:subsets}
In the cases \eqref{dir:eq:subsetcases}, every rational prime has a place of $K$ above
it of odd residue degree at most $m$, including the primes ramified in $N$, and the
bound of Proposition~\ref{dir:prop:isogenygeneral} holds with this $m$, independently of the
discriminant of $K$.
\end{corollary}

\begin{proof}
Combine Theorem~\ref{dir:thm:subsets}, Lemma~\ref{dir:lem:oddresidue} and Proposition~\ref{dir:prop:isogenygeneral}.
\end{proof}

\begin{center}
\begin{tabular}{rrrrl}
\toprule
$n$ & $k$ & $[K:\Q]$ & $m$ & cycle type attaining $m$\\
\midrule
4 & 2 & 6 & 3 & $3+1$\\
5 & 2 & 10 & 5 & $5$\\
8 & 4 & 70 & 7 & $7+1$\\
9 & 4 & 126 & 9 & $9$\\
16 & 8 & 12\,870 & 15 & $15+1$\\
17 & 8 & 24\,310 & 17 & $17$\\
\bottomrule
\end{tabular}
\end{center}
For example, a field of degree $70$ in the table admits the same universal
odd-residue-degree test as a field of degree seven. The script
\artifact{grh-arithmetic-transfers/checks/alternating.py} enumerates cycle partitions
through $n=34$ and independently visits subset orbits of individual permutations
through $n=7$.

\subsection{Heavenly abelian varieties and uniform Rasmussen--Tamagawa bounds}
\label{dir:sec:heavenly}
\sources{dirichlet-grh-research; grh-arithmetic-transfers}

Let $A$ be an abelian variety over a number field $K$ and $\ell$ a rational prime. Call
$A$ \emph{heavenly at $\ell$} if $K(A[\ell^\infty])$ lies in the maximal pro-$\ell$
extension of $K(\mu_{\ell^\infty})$ unramified outside $\ell$. Equivalently, $A$ has
good reduction away from $\ell$ and $K(A[\ell])/K(\mu_\ell)$ has $\ell$-power degree
(for elliptic curves see \cite[Lemma~2.2]{OHara2026Heavenly}). Write $\mathcal A(K,g,\ell)$ for the set of
$K$-isomorphism classes of such varieties of dimension $g$. The uniform Conjecture~2
of Rasmussen and Tamagawa \cite{RasmussenTamagawa2017Arithmetic} asks for
$\mathcal A(K,g,\ell)=\varnothing$ for $\ell>B(g,d)$ over all fields $K$ of degree $d$.
Their Theorem~5.1 assumes suitable Dedekind zeta hypotheses and treats each fixed $K$;
their Theorem~5.2 proves uniformity in odd degree under GRH. (The numbering is that of
the author manuscript, arXiv:1302.1477v1.) Lombardo established uniformity in any
degree for CM abelian varieties \cite{Lombardo2018Uniform}, and McLeman and Rasmussen
treat quadratic fields \cite{McLemanRasmussen2026Heavenly}.

\begin{lemma}[The Rasmussen--Tamagawa obstruction; \stU]\label{dir:lem:rt}
For fixed $g,d$ there exist an effective threshold $C_7(g,d)$ and a finite effective
set of integers $\mathcal M(g,d)$ with the following property. If $[K:\Q]=d$ and $A/K$
of dimension $g$ is heavenly at $\ell>C_7(g,d)$, then there is $m\in\mathcal M(g,d)$
dividing $\ell-1$ such that for every place $v\mid p\ne\ell$ of $K$,
\[
 f_v\text{ odd},\quad p^{f_v}<\ell/(4g)
 \qquad\Longrightarrow\qquad
 p^{f_v}\bmod\ell\notin(\F_\ell^\times)^m .
\]
No condition on the ramification of $p$ or $\ell$ in $K$ is needed.
\end{lemma}

\begin{proof}
Use \cite[Corollary~3.2, (3.4), Lemma~3.5, \S4.1, Proposition~4.2]{RasmussenTamagawa2017Arithmetic}.
Their hypothesis (A1) is heavenliness, and the constant $C_7$ displayed in (3.4)
suffices for (A2) uniformly in the degree. If $e$ is the gcd of their local indices
$e_\lambda$, then $m=m_\Q\mid e/2$ and $e\mid M'(2g)d$, so the positive divisors of
$M'(2g)d$ give a finite set. Here $M'(a)$ is the integer of their Lemma~3.1, the gcd
of the orders of $\GL_a(\F_{\ell'})$ over the odd primes $\ell'\ne\ell$; their
valuation formula attains each minimum at infinitely many odd primes, so removing
$\ell$ preserves the gcd. For $p\ne\ell$ the cyclotomic character at an arithmetic
Frobenius at $v$ is $p^{f_v}\bmod\ell$.

To see the sign of the obstruction, let $\chi^{i_1},\dots,\chi^{i_{2g}}$ be the
diagonal characters of $A[\ell]$. Triviality of $p^{f_v}$ modulo $m$th powers would
make all of the residual characters $\chi^{i_j+i_k-1}$ trivial on this Frobenius. If
$Q=p^{f_v}$ and $a_j$ is the integral trace of the $j$th power of Frobenius, this gives
$a_2\equiv2gQ\pmod\ell$. Weil's bound $\abs{a_2}\le2gQ$ and $4gQ<\ell$ force equality.
Every Frobenius eigenvalue is therefore $+\sqrt Q$ or $-\sqrt Q$. Since $Q$ is a
nonsquare and $a_1$ an integer, $a_1=0$. The same residual characters give
$a_1^2\equiv4g^2Q\pmod\ell$, a contradiction.
\end{proof}

\begin{theorem}[Uniform transfer; \stG]\label{dir:thm:heavenlygeneral}
Fix $g,d,r\ge1$. There is an effective bound $B(g,d,r)$ such that
$\mathcal A(K,g,\ell)=\varnothing$ for every $\ell>B(g,d,r)$ and every field $K$ of
degree $d$ satisfying the cycle hypothesis of Lemma~\ref{dir:lem:oddresidue} with this $r$.
\end{theorem}

\begin{proof}
Take the maximum of $C_7(g,d)$ and the constants $C_6(m,g,r,\Q)$ of
\cite[Proposition~2.7]{RasmussenTamagawa2017Arithmetic} over $m\in\mathcal M(g,d)$.
For $m\mid\ell-1$ let $L_{m,\ell}$ be the unique subfield of degree $m$ of
$\Q(\zeta_\ell)$. The only zeta hypothesis of that proposition here is the Riemann
hypothesis for $\zeta_{L_{m,\ell}}$, which holds by Proposition~\ref{dir:prop:abelian}. Apply the
proposition to the identity automorphism. It yields a rational prime $p\ne\ell$ with
\[
 p<(\ell/(4g))^{1/r},\qquad p\bmod\ell\in(\F_\ell^\times)^m ;
\]
the residue condition expresses complete splitting in $L_{m,\ell}$, and for $m=1$ one
may simply take $p=2$ for large $\ell$. Lemma~\ref{dir:lem:oddresidue} supplies $v\mid p$
with odd $f_v\le r$, so $p^{f_v}<\ell/(4g)$ and $p^{f_v}$ is still an $m$th power
modulo $\ell$. This contradicts Lemma~\ref{dir:lem:rt}. All fields used for the analytic
estimate are cyclotomic subfields of bounded degree.
\end{proof}

\begin{theorem}[The two sextic families; \stG]\label{dir:thm:heavenlysextic}
For every positive integer $g$ there is an effective $B(g)$ such that
$\mathcal A(K,g,\ell)=\varnothing$ whenever $[K:\Q]=6$, the normal closure of $K$ has
group $A_4$ or $A_5$, and $\ell>B(g)$. The bound is independent of the discriminant of
$K$, of the variety and of its endomorphisms. One may take $B(1)=4027$; more precisely,
a prime at which some elliptic curve over such a field is heavenly belongs to the list
\begin{equation}\label{dir:eq:heavenlylist}
\begin{gathered}
 2,3,5,7,11,13,19,23,29,31,37,43,59,67,83,107,139,163,211,\\
 283,307,331,379,499,547,643,883,907,4027 .
\end{gathered}
\end{equation}
\end{theorem}

\begin{proof}
By Lemma~\ref{dir:lem:sextic} the cycle hypothesis holds with $r=3$ for both families, so
the first assertion is Theorem~\ref{dir:thm:heavenlygeneral}. The rest of the proof concerns
$g=1$.

\emph{Inertia weights.} Let $E/K$ be a heavenly elliptic curve in either family, and
first take $\ell>5$. Its mod-$\ell$ representation on $G_K$ has semisimple characters
$\chi^i,\chi^{1-i}$ for some integer $0\le i<\ell-1$, by
\cite[Lemma~3.3]{RasmussenTamagawa2017Arithmetic} and the Weil pairing; these exponents
describe the characters on $G_K$ even if the image of $\chi|_{G_K}$ is proper. Since
$\ell\nmid\abs G$, the inertia group above $\ell$ in the normal closure is tame and
cyclic. The lengths of the cycles of an inertia generator on the embeddings equal the
ramification indices on its inertia orbits. The cycle table of Lemma~\ref{dir:lem:sextic}
therefore gives a place $\lambda\mid\ell$ of $K$ with $e(\lambda/\ell)=1$ or $3$. This
use of inertia at the torsion prime is separate from the Frobenius argument at
auxiliary primes.

At any $\lambda\mid\ell$ an elliptic curve acquires semistable reduction over an
extension of ramification index at most six. For completeness, take a short integral
equation $y^2=x^3+ax+b$ over the local field, with discriminant $\Delta$ and
$v(0)=+\infty$. If $v(j)\ge0$, then $s=v(\Delta)\le3v(a)$. From
$\Delta=-16(4a^3+27b^2)$ it follows that $s=3v(a)$ or $s=2v(b)$, since $\ell>3$; also
$s\le2v(b)$. Scale the equation by a parameter of valuation $s/12$, which exists in a
totally ramified extension of degree $12/\gcd(s,12)\le6$. The new coefficients are
integral and the new discriminant is a unit. If $v(j)<0$, then
$3v(a)=2v(b)<v(\Delta)$. Write $v(a)=2t$ and $v(b)=3t$. Scaling by a parameter of
valuation $t/2$ yields integral unit coefficients, unit $c_4$ and discriminant of
positive valuation, and the resulting reduction is multiplicative. The case $j=0$
belongs to the first alternative. In the good-reduction case the index just exhibited
is a divisor of $12$ other than $12$; in the multiplicative case it divides two. The
minimal semistability index divides any successful index
\cite[\S3.1]{RasmussenTamagawa2017Arithmetic}. At the chosen place we consequently
have
\[
 e=e_{E,K_\lambda}\,e(\lambda/\ell)\in\set{1,2,3,4,6,9,12,18}.
\]
The Tate--Oort relation \cite[(3.1)]{RasmussenTamagawa2017Arithmetic}, valid before
their large-prime assumption (A2), supplies weights $0\le j,j'\le e$ with
$ei\equiv j$ and $e(1-i)\equiv j'\pmod{\ell-1}$. For $\ell>211$ their sum must equal
$e$, since $2e<\ell-1$. Interchanging the two characters permits $j\le e/2$.

\emph{A finite test that forces equal inertia weights.} O'Hara uses powers of auxiliary
Frobenius traces to test local weights over cubic fields
\cite[\S2.3]{OHara2026Heavenly}. We use the same trace identity, retaining both
possible odd residue degrees and combining the products below by a gcd. Fix a prime
$p\in\set{2,3,5,7,11,13}$. By Lemmas~\ref{dir:lem:oddresidue} and~\ref{dir:lem:sextic} the residue
field at some $v\mid p$ has size $Q=p$ or $p^3$. At this place $E$ has good reduction.
Let $a$ be its Frobenius trace, with $a^2\le4Q$. The traces of powers satisfy
\[
 S_0=2,\quad S_1=a,\quad S_k=aS_{k-1}-QS_{k-2}.
\]
Reduction of the two eigenvalues modulo $\ell$ and the two inertia congruences give
\[
 S_e(a,Q)\equiv Q^j+Q^{e-j}\pmod\ell .
\]
For $j<e/2$ the positive integer on the right is strictly larger than $2Q^{e/2}$, by
the strict inequality of arithmetic and geometric means, whereas
$\abs{S_e}\le2Q^{e/2}$. Define
\[
 D_{e,j}=\gcd_{p\in\set{2,3,5,7,11,13}}\ \prod_{Q\in\set{p,p^3}}\
 \prod_{a^2\le4Q}\bigl(Q^j+Q^{e-j}-S_e(a,Q)\bigr),
\]
where $a$ runs over the rational integers. For unbalanced weights $\ell$ must divide
this positive gcd. The largest prime divisor of $D_{e,j}$ over $0\le j<e/2$, for each
possible $e$, is
\[
 \begin{array}{c|rrrrrrrr}
 e&1&2&3&4&6&9&12&18\\ \hline
 \text{largest prime}&13&13&109&37&109&211&109&211 .
 \end{array}
\]
The exact checker multiplies the positive terms and uses Euclid's algorithm; it divides each gcd by the
primes up to $211$ to exhaustion and checks that the remainder equals one. This proves
$j=e/2$ for every candidate $\ell>211$.

\emph{What balance forces at a good prime.} Now $(e/2)(2i-1)\equiv0\pmod{\ell-1}$.
Hence $4\mid e$, as $2i-1$ is odd, leaving $e=4$ or $12$. It also follows that
$\ell\equiv3\pmod4$ and
\[
 6(2i-1)\equiv0\pmod{\ell-1}.
\]
If $p\bmod\ell$ is a nonzero square, so is $Q=p^f$. Therefore $u=Q^{2i-1}$ has order
dividing three: write $Q=b^2$ and use the last congruence. The residual trace square
is
\[
 a^2\equiv Q(2+u+u^{-1})\equiv Q\text{ or }4Q\pmod\ell .
\]
If $4p^3<\ell$, then $0\le a^2\le4Q<\ell$, forcing equality over the integers. Both
$Q$ and $4Q$ are nonsquares, since $f=1$ or $3$. Consequently a candidate $\ell>211$
has \emph{no} prime $p$ with $4p^3<\ell$ that is a quadratic residue modulo $\ell$.

\emph{Covering all large torsion primes.} For $\ell\equiv3\pmod4$ and a prime $p$,
$\legendre p\ell=1$ exactly when $p$ splits in $\Q(\sqrt{-\ell})$. By
Theorem~\ref{dir:thm:strictcc}, every prime $\ell\equiv3\pmod4$ with $\ell>253507$ has a prime
$p$ with $4p^3<\ell$ and $\legendre p\ell=1$, so no such $\ell$ is a candidate. Direct
trial of the condition $4p^3<\ell$ through $253507$ leaves $52$ candidates larger than
$211$.

\emph{Keeping the same diagonal character at every prime.} For each remaining
$\ell>211$ retain the exponents $0\le i<\ell-1$ with $6(2i-1)\equiv0\pmod{\ell-1}$.
For $\ell\le211$ begin with all exponents, imposing only for $\ell>5$ the necessary
condition that $ei\bmod(\ell-1)\le e$ for at least one admissible $e$. For each prime
$p\ne\ell$, discard $i$ if for \emph{both} $Q=p$ and $Q=p^3$ the least absolute residue
of $Q^i+Q^{1-i}$ modulo $\ell$ exceeds $\floor{2\sqrt Q}$. For an actual curve its
fixed exponent must survive every prime. One may stop as soon as
$\floor{2\sqrt{p^3}}\ge\floor{\ell/2}$, since the option $Q=p^3$ alone allows every
residue from then on. The checker yields the list \eqref{dir:eq:heavenlylist}.
\end{proof}
\dirfiles{\artifact{dirichlet-grh-research/heavenly/check_elliptic.py},
\artifact{dirichlet-grh-research/heavenly/check_elliptic.txt}}

The tests in the proof permit the locally chosen places to vary with $p$ and impose
only necessary congruences. The large-prime step can also be obtained without the
integer certificate of Theorem~\ref{dir:thm:strictcc}: Bach and Sorenson
\cite[Theorem~5.1]{Bach1996Explicit} bound the least prime splitting in
$\Q(\sqrt{-\ell})$ by $(4\log\ell+10)^2$ under the Riemann hypothesis for its zeta
function. Put $H=3\cdot10^{13}$. For $y\ge\log H$ the function $4(4y+10)^6e^{-y}$
decreases, because $24/(4y+10)<1$, and at $y=\log H<32$ its numerator is below
$4\cdot138^6<H$; thus every $\ell>H$ has a prime $p$ with $4p^3<\ell$ and
$\legendre p\ell=1$, and the range from $253507$ to $H$ is covered by the second tail
computation of Remark~\ref{dir:rem:strict}.

\begin{corollary}[Cubic fields; \stG]\label{dir:cor:heavenlycubic}
A prime at which some elliptic curve over some cubic number field is heavenly belongs
to the list \eqref{dir:eq:heavenlylist}. In particular $25$ primes remain in the range
$\ell\ge11$.
\end{corollary}

\begin{proof}
At any rational prime the identity $\sum e_vf_v=3$ guarantees a place with odd
$f_v\in\set{1,3}$. If $\ell>5$, inertia above $\ell$ in the normal closure is cyclic,
since its group embeds in $S_3$, and its action on three points has an orbit of size
one or three. Thus the two local inputs to every numerical test in the proof of
Theorem~\ref{dir:thm:heavenlysextic} are identical.
\end{proof}

This is a numerical version of the uniform finiteness for elliptic curves over cubic
fields in \cite[Proposition~7.5]{RasmussenTamagawa2017Arithmetic}; for related balance
criteria and sieves in this family see \cite{OHara2026Heavenly}.

\begin{proposition}[Further permutation actions; \stG]\label{dir:prop:heavenlypsl}
Let $q\equiv1\pmod4$ be an odd prime power. In the natural action of $\PSL_2(\F_q)$ on
the $q+1$ points of $\mathbb P^1(\F_q)$ every element has an odd cycle of length at
most $(q+1)/2$. Hence for fixed $g$ and $q$ the conclusion of
Theorem~\ref{dir:thm:heavenlygeneral} holds uniformly over every number field realizing this
action. Moreover, if a family of fields has an odd absolute residue degree at most $r$
above every rational prime, then the extensions of fixed odd relative degree $n$ of
its members have one at most $rn$, and Theorem~\ref{dir:thm:heavenlygeneral} applies
uniformly to them.
\end{proposition}

\begin{proof}
An element with an eigenline in $\mathbb P^1(\F_q)$ fixes a point. Otherwise a lift of
determinant one has eigenvalues $\alpha,\alpha^q=\alpha^{-1}$ in $\F_{q^2}$. Therefore
$\alpha^{(q+1)/2}=\pm1$, and the projective element has order dividing the odd integer
$(q+1)/2$; every one of its cycles then has odd length at most $(q+1)/2$. The statement
on towers follows from the identity $\sum_{w\mid v}e_wf_w=n$, which gives an odd
relative residue degree at most $n$ above any place $v$.
\end{proof}

The case $q=5$ recovers the $A_5$ sextic family.

\paragraph{O'Hara's finiteness questions and descent.}
For elliptic curves we say that a pair $(E/K,\ell)$ is \emph{primitive over $K$} if no
proper intermediate field $\Q\subseteq k\subsetneq K$ supports an elliptic curve
heavenly at $\ell$ and geometrically isomorphic to $E$. In prime degree this is
O'Hara's exclusion of the classes already heavenly over $\Q$. We work inside a fixed
algebraic closure of $\Q$, and $[E]_K$ denotes a $K$-isomorphism class tagged with the
field $K$.

\begin{theorem}[O'Hara's Questions A and B; \stU]\label{dir:thm:ohara}
For cyclic cubic fields $K/\Q$ and primes $\ell\ge11$, there are finitely many pairs
$([E]_K,\ell)$ for which $E/K$ is heavenly at $\ell$ and its geometric isomorphism
class has no heavenly representative over $\Q$ at the same prime. The set of pairs
$(\ell,K)$ occurring in this way is finite.
\end{theorem}

\begin{proof}
Rasmussen and Tamagawa
\cite[Conjecture~2 and Proposition~7.5]{RasmussenTamagawa2017Arithmetic} give one bound
$N(1,3)$ valid for all cubic fields simultaneously, so $11\le\ell\le N(1,3)$. For each
such prime, McLeman and Rasmussen
\cite[Theorem~1.2, equivalently Theorem~3.1]{McLemanRasmussen2026Heavenly} with $d=3$
and $K_0=\Q$ give finitely many geometric isomorphism classes; their condition
$\ell-1>2d$ holds already at $11$. For each geometric class and prime that occur after
the exclusion, O'Hara's Lemma~5.1 excludes $j=0,1728$, and Proposition~5.2 of the same
paper gives a finite fibre of field-tagged $K$-isomorphism classes; its hypotheses hold with $p=3$
and $B=11$ \cite{OHara2026Heavenly}. A finite union proves the first statement
(Question~A), and taking the image $([E]_K,\ell)\mapsto(\ell,K)$ proves the second
(Question~B).
\end{proof}

The horizontal bound used here depends only on the degree and the dimension, and its
uniformity in $K$ permits the finite union. Corollary~\ref{dir:cor:heavenlycubic} gives $25$
candidate primes between $11$ and $4027$.

\begin{lemma}[Descent inside a fixed division field; \stU]\label{dir:lem:descent}
Suppose $E/K$ is heavenly at a prime $\ell>\max\set{2,[K:\Q]}$ and $j=j(E)\ne0,1728$.
Put $k=\Q(j)$, choose any model $A_0/k$ with invariant $j$, and set $M=k(A_0[\ell])$,
$J=K\cap M$. There exists a quadratic twist $A/J$ of $A_0$ (allowing the trivial twist)
which is heavenly at $\ell$ over $J$ and satisfies $A_K\simeq_KE$. In particular, if
the pair $(E/K,\ell)$ is primitive over $K$, then $K\subseteq M$.
\end{lemma}

\begin{proof}
\sloppy
Write $G_F$ for the absolute Galois group of a field $F$ and $\chi$ for the mod-$\ell$
cyclotomic character. The mod-$\ell$ representation of $E$ on $G_K$ is upper
triangular with diagonal $\chi^{i_1},\chi^{i_2}$
\cite[Proposition~2.4]{OHara2026Heavenly}. Since $j\ne0,1728$, there is a character
$\lambda\colon G_K\to\set{\pm1}$ giving the twist from $(A_0)_K$ to $E$. In
corresponding bases the mod-$\ell$ representation of $A_0$ on $G_K$ is upper
triangular with diagonal $\lambda\chi^{i_1},\lambda\chi^{i_2}$.

The images of $G_K$ and $G_J$ on $A_0[\ell]$ coincide: restriction identifies
$\Gal(KM/K)$ with $\Gal(M/J)$, since $M/k$ is Galois. The same basis therefore
triangularizes $G_J$. Let $\phi_1,\phi_2$ be its diagonal characters. The determinant
$\chi$ factors through $\Gal(M/J)$ by the Weil pairing. Thus the characters
$\phi_1\chi^{-i_1}$ and $\phi_2\chi^{-i_2}$ on $G_J$ factor through that shared finite
image, and both restrict to $\lambda$ on $G_K$. They agree and take values in
$\set{\pm1}$ everywhere. Twisting $(A_0)_J$ by this common character gives a curve $A$
which becomes $E$ over $K$ and whose diagonal characters on $G_J$ are the two
cyclotomic powers. In particular $J(A[\ell])/J(\mu_\ell)$ has $\ell$-power degree.

For each $s\ge1$ put $F_s=J(A[\ell^s])$. The congruence kernels in
$\GL_2(\Z/\ell^s)$ are $\ell$-groups, so $F_s/J(\mu_\ell)$ has $\ell$-power degree. At
every finite prime of $J$ away from $\ell$, the inertia order of $F_s/J$ is
consequently a power of $\ell$. Since $KF_s/K$ is unramified there, multiplicativity of
ramification indices in the two towers shows that this inertia order divides the
corresponding ramification index in $K/J$. It is at most $[K:J]<\ell$ and hence equals
one. The criterion of N\'eron--Ogg--Shafarevich gives good reduction of $A$ away from
$\ell$. This proves the descent, and primitivity then forces $K=J$.
\end{proof}

\begin{theorem}[Odd degree; \stG]\label{dir:thm:heavenlyodd}
Fix an odd integer $n>1$ and let $K$ range over all number fields of degree $n$. There
are finitely many field-tagged pairs $([E]_K,\ell)$ primitive over $K$ with $E$
heavenly at a prime $\ell>2n+1$. For an odd prime $n$ the exclusion of the classes
heavenly over $\Q$ alone suffices. The same assertions hold upon restricting to Galois
fields.
\end{theorem}

\begin{proof}
In the proof of \cite[Theorem~5.3]{RasmussenTamagawa2017Arithmetic} take the fixed base
$F=\Q$. The Chebotarev application there uses the field
$\widetilde L=\Q(\mu_\ell)_{m_\Q}\widetilde F$; here $\widetilde F=\Q$ and the
subscript indicates a subfield of the cyclotomic field. Its Dedekind zeta function
factors into Dirichlet $L$-functions, so the GRH input for this application is
available by Proposition~\ref{dir:prop:abelian}. The choice of the auxiliary prime is uniform in
the field $K$ of degree $n$. The passage to $K$ uses the existence of a place of odd
residue degree at most $n$ over that prime, which follows from $\sum e_vf_v=n$ being
odd. The rest of the cited proof gives a bound $\ell\le N(1,n)$ uniform in $K$; this
specializes the proof of their more general theorem under Dedekind GRH to the rational
base.

For each remaining $\ell$, \cite[Theorem~3.1]{McLemanRasmussen2026Heavenly} gives
finitely many $j$. The exceptional invariants $0,1728$ are impossible in odd degree
for $\ell>3$. Indeed the CM-field containment used in
\cite[Lemma~5.1]{OHara2026Heavenly} forces $K\Q(\sqrt{-3})$ or $K\Q(i)$, respectively,
into $K(E[\ell])$. The latter extension over $K(\mu_\ell)$ has odd degree, so the
corresponding quadratic field would lie in $K(\mu_\ell)$. It is disjoint from
$\Q(\mu_\ell)$ because of ramification at $3$ or $2$, whereas
$[K(\mu_\ell):\Q(\mu_\ell)]$ divides the odd degree $n$. This argument uses no
normality of $K$. For each other $j$, primitivity and Lemma~\ref{dir:lem:descent} force $K$
to lie in the fixed finite extension $k(A_0[\ell])$, which allows only finitely many
$K$. For each of them Shafarevich's finiteness theorem for elliptic curves with good
reduction away from $\ell$ gives finitely many $K$-isomorphism classes. Taking the
unions finishes the proof.
\end{proof}

\begin{theorem}[The two sextic families; \stG]\label{dir:thm:heavenlysexticfinite}
Let $K$ range through all sextic fields with normal closure of group $A_5$. There are
finitely many pairs $([E]_K,\ell)$ with $\ell\ge17$, $E/K$ heavenly at $\ell$, and no
heavenly representative of the geometric class over $\Q$ at $\ell$. For sextic fields
with normal closure of group $A_4$ the same finiteness holds when one excludes
representatives over every proper intermediate field. In each case the pairs
$(\ell,K)$ form a finite set.
\end{theorem}

\begin{proof}
By Theorem~\ref{dir:thm:heavenlysextic}, $\ell\le4027$ for these fields. For each such $\ell$
the theorem of McLeman and Rasmussen applies with $d=6$, since $\ell-1>12$. It remains
to treat the exceptional invariants and the proper-subfield condition.

Let $N$ be the normal closure, $G=\Gal(N/\Q)$ and $F=\Q(\mu_\ell)$. If $j=0$ or
$1728$, the CM-field containment in O'Hara's proof puts $D=\Q(\sqrt{-3})$ or $\Q(i)$
in $K(E[\ell])$. Over $KF$ this division field has odd degree, so
$D\subseteq KF\subseteq NF$. We rule this out directly. The group
$\Gal(NF/\Q)$ lies in $G\times\Gal(F/\Q)$ and surjects onto $G$. Its commutator
subgroup is $G'\times\set{1}$: the second factor is abelian and every commutator in
$G$ lifts. The largest abelian subfield of $NF$ is therefore $N^{G'}F$. Its degree
over $F$ divides three for $A_4$ and equals one for $A_5$. Since $D$ is quadratic and
ramifies at $2$ or $3$, while $F$ ramifies only at $\ell\ge17$, the quadratic extension
$DF/F$ cannot lie in that abelian subfield. This excludes both exceptional invariants.

Lemma~\ref{dir:lem:descent} now bounds the fields for each $j,\ell$ under the
proper-subfield exclusion, and Shafarevich's theorem bounds the curves for each field.
In the $A_5$ action the point stabilizer has order ten; a strictly larger proper
subgroup would have index two or three, impossible by simplicity of $A_5$ and the
associated coset action. Hence $K$ has no proper intermediate field except $\Q$. In
the $A_4$ action the stabilizer of order two lies in the Klein four subgroup, giving a
cubic intermediate field, which the stated exclusion takes into account.
\end{proof}
